% ZXTune_LHUNA — manuel de référence
% Compilation : lualatex lhuna-bible ; lualatex lhuna-bible (deux passes pour la table des matières)
% Les tables de formats sont générées : voir tools/formats.py et gen/formats.tex
\documentclass[11pt,a4paper,twoside,openright]{book}

\usepackage{fontspec}
\usepackage[french]{babel}
\usepackage[margin=1.5cm,bindingoffset=0.5cm]{geometry}
\usepackage{graphicx}
\usepackage{longtable}
\usepackage{booktabs}
\usepackage{array}
\usepackage{xcolor}
\usepackage[most]{tcolorbox}
\usepackage{fancyhdr}
\usepackage{titlesec}
\usepackage{enumitem}
\usepackage[hidelinks]{hyperref}
\usepackage{xspace}
% nom du logiciel ; la marque et le logo restent « LHUNA »
\newcommand{\soft}{ZXTune\_LHUNA\xspace}

\graphicspath{{images/}}

\definecolor{lhunacyan}{RGB}{0,229,255}
\definecolor{lhunanight}{RGB}{11,14,20}
\definecolor{lhunagrey}{RGB}{110,118,130}

\setlist{itemsep=2pt,topsep=4pt}
\titleformat{\chapter}[display]{\normalfont\huge\bfseries}{\color{lhunacyan}\chaptertitlename\ \thechapter}{12pt}{\Huge}
\titleformat{\section}{\normalfont\Large\bfseries}{\color{lhunacyan}\thesection}{1em}{}

\pagestyle{fancy}
\fancyhf{}
\fancyhead[LE]{\small\leftmark}
\fancyhead[RO]{\small\rightmark}
\fancyfoot[C]{\small\thepage}
\renewcommand{\headrulewidth}{0.3pt}

% Encadré pour non-techniciens
\newtcolorbox{enclair}{colback=lhunacyan!6,colframe=lhunacyan!70!black,title=En clair,fonttitle=\bfseries,breakable}
% Encadré technique
\newtcolorbox{technique}{colback=lhunagrey!8,colframe=lhunagrey,title=Technique,fonttitle=\bfseries,breakable}
% Partie restant à rédiger (repérable par recherche de « À RÉDIGER »)
\newcommand{\arediger}[1]{\begin{tcolorbox}[colback=red!5,colframe=red!60!black,title=À RÉDIGER,breakable]#1\end{tcolorbox}}

% Capture d'écran centrée avec légende ; largeur relative à la ligne
\newcommand{\capture}[3][0.45]{%
  \begin{figure}[htbp]\centering
  \includegraphics[width=#1\linewidth]{#2}
  \caption{#3}
  \end{figure}}

% Deux captures côte à côte
\newcommand{\capturedouble}[4]{%
  \begin{figure}[htbp]\centering
  \includegraphics[width=0.45\linewidth]{#1}\hfill
  \includegraphics[width=0.45\linewidth]{#2}
  \caption{#3 (gauche) et #4 (droite)}
  \end{figure}}

% generated by tools/formats.py, do not edit

\newcommand{\lhunaModulesCount}{287}
\newcommand{\lhunaContainersCount}{97}
\newcommand{\lhunaFormatsTable}{%
\begin{longtable}{@{}l p{0.46\textwidth} l l@{}}
\toprule \textbf{Id} & \textbf{Format} & \textbf{Type} & \textbf{Famille} \\ \midrule \endhead
AS0 & ASC Sound Master v0.x & Tracker module & AY-3-8910/YM2149 \\
ASC & ASC Sound Master v1.x & Tracker module & AY-3-8910/YM2149 \\
AYC & CPC AYC & Register stream & AY-3-8910/YM2149 \\
FTC & Fast Tracker v1.x & Tracker module & AY-3-8910/YM2149 \\
GTR & Global Tracker v1.x & Tracker module & AY-3-8910/YM2149 \\
PSC & Pro Sound Creator v1.xx & Tracker module & AY-3-8910/YM2149 \\
PSG & Programmable Sound Generator & Register stream & AY-3-8910/YM2149 \\
PSM & Pro Sound Maker & Tracker module & AY-3-8910/YM2149 \\
PT1 & Pro Tracker v1.x & Tracker module & AY-3-8910/YM2149 \\
PT2 & Pro Tracker v2.x & Tracker module & AY-3-8910/YM2149 \\
SQT & SQ-Tracker Compiled & Tracker module & AY-3-8910/YM2149 \\
ST1 & Sound Tracker v1.x & Tracker module & AY-3-8910/YM2149 \\
ST3 & Sound Tracker v3.x Compiled & Tracker module & AY-3-8910/YM2149 \\
STC & Sound Tracker v1.x Compiled & Tracker module & AY-3-8910/YM2149 \\
STP & Sound Tracker Pro Compiled & Tracker module & AY-3-8910/YM2149 \\
TXT & VortexTracker II & Tracker module & AY-3-8910/YM2149 \\
VTX & VTX (Vortex Project) & Register stream & AY-3-8910/YM2149 \\
YM & YM (ST-Sound Project) & Register stream & AY-3-8910/YM2149 \\
YM & YM (ST-Sound Project) Packed & Register stream & AY-3-8910/YM2149 \\
AY & AY/EMUL & Memory dump & AY-3-8910/YM2149, Beeper \\
PT3 & Pro Tracker v3.x & Tracker module & AY-3-8910/YM2149, TurboSound \\
RMT & Raster Music Tracker & Memory dump & Atari POKEY \\
SAP & Slight Atari Player Sound Format & Memory dump & Atari POKEY \\
DSF & Dreamcast Sound Format & Memory dump & Console: Dreamcast \\
GSF & GameBoy Advance Sound Format & Memory dump & Console: Game Boy Advance \\
USF & Ultra64 Sound Format & Memory dump & Console: Nintendo 64 \\
2SF & Nintendo DS Sound Format & Memory dump & Console: Nintendo DS \\
NCSF & Nitro Composer Sound Format & Memory dump & Console: Nintendo DS \\
PSF & Playstation Sound Format & Memory dump & Console: PlayStation \\
PSF2 & Playstation2 Sound Format & Memory dump & Console: PlayStation 2 \\
SSF & Sega Saturn Sound Format & Memory dump & Console: Sega Saturn \\
SNSF & Super NES Sound Format & Memory dump & Console: Super Nintendo \\
GBS & GameBoy Sound & Memory dump & Game Boy \\
2DX9 & Konami 2DX9 & Register stream & Game audio stream \\
2MSF & Banpresto RIFF wrapper & Register stream & Game audio stream \\
AAC & tri-Ace ASKA engine & Register stream & Game audio stream \\
AAX & CRI AAX & Register stream & Game audio stream \\
ACB & CRI Atom Cursheet Binary & Register stream & Game audio stream \\
ACM & InterPlay ACM & Register stream & Game audio stream \\
ADP & Nintendo ADP (raw) & Register stream & Game audio stream \\
ADP & Nex Entertainment NXAP & Register stream & Game audio stream \\
ADP & Xilam DERF & Register stream & Game audio stream \\
ADS & Sony AudioStream & Register stream & Game audio stream \\
ADS & Midway ADS & Register stream & Game audio stream \\
ADX & CRI ADX & Register stream & Game audio stream \\
AHX & CRI AHX & Register stream & Game audio stream \\
AIX & CRI AIX & Register stream & Game audio stream \\
AKB & Square Enix AKB & Register stream & Game audio stream \\
AKB & Square Enix AKB2 & Register stream & Game audio stream \\
AST & Nintendo AST & Register stream & Game audio stream \\
AUD & Westwood Studios AUD & Register stream & Game audio stream \\
AUS & Capcom AUS & Register stream & Game audio stream \\
AWB & CRI Atom Wave Bank & Register stream & Game audio stream \\
AWC & Rockstar AWC & Register stream & Game audio stream \\
BAF & Bizzare Creations bank file & Register stream & Game audio stream \\
BAF & Bizzare Creations band file (bad rips) & Register stream & Game audio stream \\
BANK & FSB5 FEV & Register stream & Game audio stream \\
BAO & Ubisoft BAO (atomic) & Register stream & Game audio stream \\
BGW & Square Enix BGW & Register stream & Game audio stream \\
BIK & RAD Game Tools Bink & Register stream & Game audio stream \\
BMP & Konami BMP & Register stream & Game audio stream \\
BNK & Sony BNK (LE) & Register stream & Game audio stream \\
BNK & Sony BNK (BE) & Register stream & Game audio stream \\
BNK & Audiokinetic Wwise Soundbank Container & Register stream & Game audio stream \\
BNK & Electronic Arts BNK & Register stream & Game audio stream \\
BWAV & Nintendo BWAV & Register stream & Game audio stream \\
CAF & tri-Crescendo CAF & Register stream & Game audio stream \\
CKD & Ubisoft CKD & Register stream & Game audio stream \\
CKS & Cricket Audio Stream & Register stream & Game audio stream \\
CPK & CRI CPK container & Register stream & Game audio stream \\
CSTM & Nintendo CSTM & Register stream & Game audio stream \\
CWAV & Nintendo CWAV & Register stream & Game audio stream \\
DSP & Nintendo ADPCM & Register stream & Game audio stream \\
DSP & Nintendo ADPCM (LE) & Register stream & Game audio stream \\
DSP & Koei Tecmo WiiBGM & Register stream & Game audio stream \\
DSP & Namco CSTR & Register stream & Game audio stream \\
DSP & Retro Studios RS03 & Register stream & Game audio stream \\
DSPW & Capcom DSPW & Register stream & Game audio stream \\
DVI & Capcom IDVI & Register stream & Game audio stream \\
FSB & FMOD Sample Bank (v1-v5) & Register stream & Game audio stream \\
FSTM & Nintendo FSTM & Register stream & Game audio stream \\
FWAV & Nintendo FWAV & Register stream & Game audio stream \\
GCA & Terminal Reality GCA & Register stream & Game audio stream \\
GENH & Generic Header & Register stream & Game audio stream \\
H4M & Hudson HVQM4 & Register stream & Game audio stream \\
HCA & CRI HCA & Register stream & Game audio stream \\
HPS & HAL Labs HALPST & Register stream & Game audio stream \\
HWAS & Vicarious Visions HWAS & Register stream & Game audio stream \\
IDSP & Namco IDSP (NUB container) & Register stream & Game audio stream \\
IDSP & Various IDSP & Register stream & Game audio stream \\
IFS & Konami Arcade Games Container & Register stream & Game audio stream \\
IKM & MiCROViSiON container & Register stream & Game audio stream \\
IMUS & LucasArts iMUSE & Register stream & Game audio stream \\
ISD & Inti Creates OGG Vorbis & Register stream & Game audio stream \\
ISWS & Sumo Digital iSWS & Register stream & Game audio stream \\
KCEY & Konami KCE Yokohama & Register stream & Game audio stream \\
KTSS & Koei Tecmo KTSS & Register stream & Game audio stream \\
KVS & Koei Tecmo KVS & Register stream & Game audio stream \\
LOGG & Encrypted OGG Vorbis & Register stream & Game audio stream \\
LOGG & Corrupted OGG Vorbis & Register stream & Game audio stream \\
LOPUS & Capcom OPUS container & Register stream & Game audio stream \\
LWAV & Ubisoft RIFF & Register stream & Game audio stream \\
LWAV & Ubisoft LyN & Register stream & Game audio stream \\
M4A & Apple QuickTime Container & Register stream & Game audio stream \\
MAB & Square Enix MAB & Register stream & Game audio stream \\
MCA & Capcom MCA & Register stream & Game audio stream \\
MIC & Sony MultiStream MIC & Register stream & Game audio stream \\
MIC & Koei Tecmo MIC & Register stream & Game audio stream \\
MSF & Sony MultiStream File & Register stream & Game audio stream \\
MSF & Konami MSF & Register stream & Game audio stream \\
MTX & Matric Software container & Register stream & Game audio stream \\
MUL & Crystal Dynamics MUL & Register stream & Game audio stream \\
MUL & Crystal Dynamics MUL (BE) & Register stream & Game audio stream \\
MUSC & Krome MUSC & Register stream & Game audio stream \\
MUSX & Eurocom MUSX & Register stream & Game audio stream \\
NAAC & Namco AAC & Register stream & Game audio stream \\
NPS & Namco NPSF & Register stream & Game audio stream \\
NSOPUS & Nihon Falcom FDK & Register stream & Game audio stream \\
NUB & Namco nuSound container & Register stream & Game audio stream \\
NUS3 & Namco nuSound container v3 & Register stream & Game audio stream \\
NWA & VisualArt's NWA & Register stream & Game audio stream \\
NXA & Entergram NXA & Register stream & Game audio stream \\
OMA & Sony OpenMG audio & Register stream & Game audio stream \\
OPUS & Nintendo OPUS & Register stream & Game audio stream \\
OPUS & Various OPUS containers & Register stream & Game audio stream \\
P3D & Radical P3D & Register stream & Game audio stream \\
PDT & Hudson Splitted Stream container & Register stream & Game audio stream \\
PDT & Hudson Stream Container & Register stream & Game audio stream \\
PONA & Policenauts BGM & Register stream & Game audio stream \\
RAK & Ubisoft RAKI & Register stream & Game audio stream \\
RAK & Ubisoft RAKI (modified) & Register stream & Game audio stream \\
RFRM & Retro Studios RFRM & Register stream & Game audio stream \\
RPGMVO & Encrypted OGG Vorbis & Register stream & Game audio stream \\
RSD & Radical RSD & Register stream & Game audio stream \\
RSD & RedSparc RSD & Register stream & Game audio stream \\
RSTM & Nintendo RSTM & Register stream & Game audio stream \\
RWAV & Nintendo RWAV & Register stream & Game audio stream \\
RWS & RenderWare Stream & Register stream & Game audio stream \\
RXW & Sony XWS (bad rips) & Register stream & Game audio stream \\
SAB & Square Enix SAB & Register stream & Game audio stream \\
SAD & Procyon Studio SADL & Register stream & Game audio stream \\
SCD & Square Enix SCD & Register stream & Game audio stream \\
SCHL & Electronic Arts SCHL & Register stream & Game audio stream \\
SD9 & Konami SD9 & Register stream & Game audio stream \\
SDF & Beyond Reality SDF & Register stream & Game audio stream \\
SEB & Game Arts SEB & Register stream & Game audio stream \\
SEG & Stormfront SEG & Register stream & Game audio stream \\
SMK & RAD Game Tools SMACKER & Register stream & Game audio stream \\
SMP & Infernal Engine SMP & Register stream & Game audio stream \\
SNR & Electronic Arts SNR & Register stream & Game audio stream \\
SPS & Electronic Arts SPS & Register stream & Game audio stream \\
SPSD & Naomi SPSD & Register stream & Game audio stream \\
STER & Alchemy STER & Register stream & Game audio stream \\
STM & Intelligent Systems STM & Register stream & Game audio stream \\
STR & Sega Stream Asset Builder STR & Register stream & Game audio stream \\
STR & 3DO STR & Register stream & Game audio stream \\
STRM & Nintendo STRM & Register stream & Game audio stream \\
SVAG & Konami SVAG & Register stream & Game audio stream \\
SVAG & SNK SVAG & Register stream & Game audio stream \\
SWAV & Nintendo SWAV & Register stream & Game audio stream \\
THP & Nintendo THP & Register stream & Game audio stream \\
VAG & Sony VAG & Register stream & Game audio stream \\
VGS & Princess Soft VGS & Register stream & Game audio stream \\
VPK & SCE America VPK & Register stream & Game audio stream \\
VXN & Gameloft VXN & Register stream & Game audio stream \\
WAVE & EngineBlack WAVE & Register stream & Game audio stream \\
WEM & Audiokinetic Wwise & Register stream & Game audio stream \\
WMSF & Banpresto MSF wrapper & Register stream & Game audio stream \\
XA & Sony XA & Register stream & Game audio stream \\
XA & Sony XA (raw) & Register stream & Game audio stream \\
XA & Maxis XA & Register stream & Game audio stream \\
XMA & Microsoft XMA & Register stream & Game audio stream \\
XNB & Microsoft XNA Game Studio XNB & Register stream & Game audio stream \\
XVAG & Sony XVAG & Register stream & Game audio stream \\
XWB & Microsoft XACT Wave Bank & Register stream & Game audio stream \\
XWC & Starbreeze XWC & Register stream & Game audio stream \\
XWM & Microsoft XWMA & Register stream & Game audio stream \\
XWS & Sony XWS & Register stream & Game audio stream \\
XWV & feelplus XWAV & Register stream & Game audio stream \\
YDSP & Yuke DSP & Register stream & Game audio stream \\
ZSD & ZSD container & Register stream & Game audio stream \\
ZSS & Z-Axis ZSND & Register stream & Game audio stream \\
HES & Home Entertainment System & Memory dump & HuC6280 \\
MIDI & Standard MIDI File & Tracker module & MIDI (Roland MT-32) \\
SID & Commodore64 SID/RSID/PSID & Memory dump & MOS6581/8580 \\
NSF & NES Sound Format & Memory dump & Ricoh 2A03/2A07 \\
NSFE & Extended Nintendo Sound Format & Memory dump & Ricoh 2A03/2A07 \\
COP & E-Tracker v1.x & Tracker module & SAA1099 \\
SPC & SNES SPC700 & Memory dump & SPC700/S-DSP \\
GYM & Genesis YM2612 (unpacked) & Register stream & Several chips \\
KSS & KSS Music Format & Memory dump & Several chips \\
KSSX & KSS Extended Music Format & Memory dump & Several chips \\
MTC & Multitrack Container & Multichip & Several chips \\
S98 & Sound 98 & Register stream & Several chips \\
VGM & Video Game Music & Register stream & Several chips \\
AAC & Advanced Audio Coding (raw ADTS) & Register stream & Standard audio \\
AC3 & Dolby AC-3 & Register stream & Standard audio \\
AIFC & Apple Compressed Audio Interchange File Format & Register stream & Standard audio \\
AIFF & Apple Audio Interchange File Format & Register stream & Standard audio \\
APE & Monkey's Audio & Register stream & Standard audio \\
ASF & Windows Media Audio & Register stream & Standard audio \\
CAF & Apple Core Audio Format File & Register stream & Standard audio \\
FLAC & Free Lossless Audio Codec & Register stream & Standard audio \\
MP3 & MPEG Audio Layer & Register stream & Standard audio \\
MPC & Musepack & Register stream & Standard audio \\
OGG & OGG Vorbis & Register stream & Standard audio \\
OPUS & OGG Opus & Register stream & Standard audio \\
TAK & Tom's lossless Audio Kompressor & Register stream & Standard audio \\
WAV & Waveform Audio File Format & Register stream & Standard audio \\
667 & Composer 667 & Tracker module & Tracker (samples) \\
669 & 669 Composer / UNIS 669 & Tracker module & Tracker (samples) \\
AHX & Abyss' Highest Experience & Tracker module & Tracker (samples) \\
AMF & ASYLUM Music Format & Tracker module & Tracker (samples) \\
AMF & DSMI / Digital Sound And Music Interface & Tracker module & Tracker (samples) \\
AMS & Extreme's Tracker & Tracker module & Tracker (samples) \\
AMS & Velvet Studio & Tracker module & Tracker (samples) \\
C67 & Composer 670 & Tracker module & Tracker (samples) \\
CBA & Chuck Biscuits / Black Artist & Tracker module & Tracker (samples) \\
CHI & Chip Tracker & Tracker module & Tracker (samples) \\
DBM & DigiBooster Pro & Tracker module & Tracker (samples) \\
DBM & Digi Booster & Tracker module & Tracker (samples) \\
DMF & X-Tracker & Tracker module & Tracker (samples) \\
DMM & Digital Music Maker v1.x & Tracker module & Tracker (samples) \\
DSM & Digital Sound Interface Kit RIFF & Tracker module & Tracker (samples) \\
DST & Digital Studio & Tracker module & Tracker (samples) \\
DSYM & Digital Symphony & Tracker module & Tracker (samples) \\
DTM & Digital Tracker / Digital Home Studio & Tracker module & Tracker (samples) \\
DTT & Desktop Tracker & Tracker module & Tracker (samples) \\
EMOD & Quadra Composer & Tracker module & Tracker (samples) \\
ET1 & Extreme Tracker v1.x & Tracker module & Tracker (samples) \\
ETX & EasyTrax & Tracker module & Tracker (samples) \\
FAR & Farandole Composer & Tracker module & Tracker (samples) \\
FC & Future Composer 1.0 - 1.4 & Tracker module & Tracker (samples) \\
FMT & Davey W Taylor's FM Tracker & Tracker module & Tracker (samples) \\
FNK & Funktracker & Tracker module & Tracker (samples) \\
FTM & Face The Music & Tracker module & Tracker (samples) \\
GDM & BWSB Soundsystem & Tracker module & Tracker (samples) \\
GMC & Game Music Creator & Tracker module & Tracker (samples) \\
GTK & Graoumf Tracker 1/2 & Tracker module & Tracker (samples) \\
HVL & Hively Tracker & Tracker module & Tracker (samples) \\
IMF & Imago Orpheus & Tracker module & Tracker (samples) \\
IMS & Images Music System & Tracker module & Tracker (samples) \\
IT & Impulse Tracker & Tracker module & Tracker (samples) \\
J2B & Galaxy Sound System & Tracker module & Tracker (samples) \\
KRIS & ChipTracker & Tracker module & Tracker (samples) \\
LIQ & Liquid Tracker & Tracker module & Tracker (samples) \\
LIQ & Liquid Tracker NO & Tracker module & Tracker (samples) \\
MDL & Digitrakker & Tracker module & Tracker (samples) \\
MED & MED 1.12 MED2 & Tracker module & Tracker (samples) \\
MED & MED 2.00 MED3 & Tracker module & Tracker (samples) \\
MED & MED 2.10 MED4 & Tracker module & Tracker (samples) \\
MED & OctaMED / MED Soundstudio & Tracker module & Tracker (samples) \\
MO3 & Un4seen MO3 & Tracker module & Tracker (samples) \\
MOD & ProTracker 3.6 & Tracker module & Tracker (samples) \\
MOD & Generic MOD-compatible & Tracker module & Tracker (samples) \\
MOD & MnemoTroN SoundTracker (MOD-compatible) & Tracker module & Tracker (samples) \\
MOD & Ice Tracker (MOD-compatible) & Tracker module & Tracker (samples) \\
MOD & Ultimate Soundtracker / etc (MOD Compatible) & Tracker module & Tracker (samples) \\
MT2 & Mad Tracker 2.xx & Tracker module & Tracker (samples) \\
MTM & MultiTracker & Tracker module & Tracker (samples) \\
OKT & Oktalyzer & Tracker module & Tracker (samples) \\
PDT & ProDigi Tracker v0.0x & Tracker module & Tracker (samples) \\
PLM & Disorder Tracker 2 & Tracker module & Tracker (samples) \\
PSM & Epic MegaGames MASI & Tracker module & Tracker (samples) \\
PSM & Epic MegaGames MASI (Old Version) & Tracker module & Tracker (samples) \\
PTM & PolyTracker & Tracker module & Tracker (samples) \\
PUMA & PumaTracker & Tracker module & Tracker (samples) \\
RTM & Real Tracker 2 & Tracker module & Tracker (samples) \\
S3M & ScreamTracker 3 & Tracker module & Tracker (samples) \\
SFX & SoundFX 1.x & Tracker module & Tracker (samples) \\
SFX & SoundFX 2.0 / MultiMedia Sound & Tracker module & Tracker (samples) \\
SQD & SQ Digital Tracker & Tracker module & Tracker (samples) \\
STIM & Slamtilt & Tracker module & Tracker (samples) \\
STM & ScreamTracker 2 & Tracker module & Tracker (samples) \\
STP & Soundtracker Pro II & Tracker module & Tracker (samples) \\
STR & Sample Tracker v2.x & Tracker module & Tracker (samples) \\
STX & STMIK 0.2 & Tracker module & Tracker (samples) \\
SYMMOD & Symphonie & Tracker module & Tracker (samples) \\
TCB & TCB Tracker & Tracker module & Tracker (samples) \\
ULT & UltraTracker & Tracker module & Tracker (samples) \\
V2M & Farbrausch V2 Synthesizer System & Tracker module & Tracker (samples) \\
XM & FastTracker II & Tracker module & Tracker (samples) \\
TF0 & TFM Music Maker v0.1-1.2 & Tracker module & TurboFM (YM2203) \\
TFC & TurboFM Compiled Dump & Register stream & TurboFM (YM2203) \\
TFD & TurboFM Dump & Register stream & TurboFM (YM2203) \\
TFE & TFM Music Maker v1.3+ & Tracker module & TurboFM (YM2203) \\
TS & TurboSound & Multichip & TurboSound \\
\bottomrule \end{longtable}}
\newcommand{\lhunaContainersTable}{%
\begin{longtable}{@{}l p{0.55\textwidth} l@{}}
\toprule \textbf{Id} & \textbf{Conteneur} & \textbf{Type} \\ \midrule \endhead
7ZIP & 7zip & Archive \\
GZIP & GZip & Archive \\
HOBETA & Hobeta & Archive \\
HRIP & Hrip (Hrust RiP archiver) & Archive \\
LHA & LHA (LHArc) & Archive \\
RAR & RAR & Archive \\
UMX & UMX (Unreal Music eXperience) & Archive \\
ZDATA & Zdata & Archive \\
ZIP & ZIP & Archive \\
ZXZIP & ZXZip & Archive \\
CC3 & CodeCruncher v3.x & Compressor \\
CC4 & CompressorCode v.4 by ZYX & Compressor \\
CC4PLUS & CompressorCode v.4+ by ZYX & Compressor \\
CHARPRES & CharPres & Compressor \\
DSQ & DataSqueezer v4.x & Compressor \\
ESV & ESV Cruncher & Compressor \\
GAM & GamePacker & Compressor \\
GAMPLUS & GamePacker+ & Compressor \\
HRUM & Hrum v3.x & Compressor \\
HRUST1 & Hrust v1.x & Compressor \\
HRUST2 & Hrust v2.1 & Compressor \\
HRUST23 & Hrust v2.3 & Compressor \\
LZH1 & LZH Compressor v1.4 & Compressor \\
LZH2 & LZH Compressor v2.4 & Compressor \\
LZS & ASC LZS Pack & Compressor \\
MEGALZ & MegaLZ & Compressor \\
MSP & MicroSpace Packer v1.x & Compressor \\
MUSE & MUSE Compressor & Compressor \\
PACK2 & Pack v2.x & Compressor \\
PCD61 & Powerfull Code Decreaser v6.1 & Compressor \\
PCD61i & Powerfull Code Decreaser v6.1i & Compressor \\
PCD62 & Powerfull Code Decreaser v6.2 & Compressor \\
TLZ & Turbo-LZ v1.x & Compressor \\
TLZP & Turbo\#LZ v1.x (internal) & Compressor \\
TRUSH & Trush Compressor v3.x & Compressor \\
DSK & DSK Image & Disk image \\
FDI & FDI (Full Disk Image) & Disk image \\
SCL & SCL (SINCLAIR) & Disk image \\
TD0 & TD0 (TeleDisk Image) & Disk image \\
TRD & TRD (TR-DOS) & Disk image \\
AAC & Multitrack tri-Ace ASKA engine & Multitrack file \\
ACB & Multitrack CRI Atom Cursheet Binary & Multitrack file \\
AKB & Multitrack Square Enix AKB2 & Multitrack file \\
ASF & Multitrack Windows Media Audio & Multitrack file \\
AWB & Multitrack CRI Atom Wave Bank & Multitrack file \\
AWC & Multitrack Rockstar AWC & Multitrack file \\
AY & Multi-AY/EMUL & Multitrack file \\
BAF & Multitrack Bizzare Creations bank file & Multitrack file \\
BAF & Multitrack Bizzare Creations band file (bad rips) & Multitrack file \\
BIK & Multitrack RAD Game Tools Bink & Multitrack file \\
BNK & Multitrack Sony BNK (LE) & Multitrack file \\
BNK & Multitrack Sony BNK (BE) & Multitrack file \\
BNK & Multitrack Audiokinetic Wwise Soundbank Container & Multitrack file \\
BNK & Multitrack Electronic Arts BNK & Multitrack file \\
CPK & Multitrack CRI CPK container & Multitrack file \\
FSB & Multitrack FMOD Sample Bank (v1-v5) & Multitrack file \\
GBS & Multitrack GameBoy Sound & Multitrack file \\
H4M & Multitrack Hudson HVQM4 & Multitrack file \\
HES & Multitrack Home Entertainment System & Multitrack file \\
IFS & Multitrack Konami Arcade Games Container & Multitrack file \\
KSSX & Multitrack KSS Extended Music Format & Multitrack file \\
M4A & Multitrack Apple QuickTime Container & Multitrack file \\
MAB & Multitrack Square Enix MAB & Multitrack file \\
MUSX & Multitrack Eurocom MUSX & Multitrack file \\
NSF & Multitrack NES Sound Format & Multitrack file \\
NSFE & Multitrack Extended Nintendo Sound Format & Multitrack file \\
NUB & Multitrack Namco nuSound container & Multitrack file \\
NUS3 & Multitrack Namco nuSound container v3 & Multitrack file \\
PDT & Multitrack Hudson Stream Container & Multitrack file \\
RWS & Multitrack RenderWare Stream & Multitrack file \\
RXW & Multitrack Sony XWS (bad rips) & Multitrack file \\
SAB & Multitrack Square Enix SAB & Multitrack file \\
SAP & Multitrack Slight Atari Player Sound Format & Multitrack file \\
SCD & Multitrack Square Enix SCD & Multitrack file \\
SID & Multitrack Commodore64 SID/RSID/PSID & Multitrack file \\
VXN & Multitrack Gameloft VXN & Multitrack file \\
XA & Multitrack Sony XA & Multitrack file \\
XA & Multitrack Sony XA (raw) & Multitrack file \\
XVAG & Multitrack Sony XVAG & Multitrack file \\
XWB & Multitrack Microsoft XACT Wave Bank & Multitrack file \\
XWM & Multitrack Microsoft XWMA & Multitrack file \\
XWS & Multitrack Sony XWS & Multitrack file \\
ZSS & Multitrack Z-Axis ZSND & Multitrack file \\
COMPILEDASC0 & ASC Sound Master v0.x player & Player decompiler \\
COMPILEDASC1 & ASC Sound Master v1.x player & Player decompiler \\
COMPILEDASC2 & ASC Sound Master v2.x player & Player decompiler \\
COMPILEDPT24 & Pro Tracker v2.40 Phantom Family player & Player decompiler \\
COMPILEDPTU13 & Pro Tracker Utility v1.3 player & Player decompiler \\
COMPILEDST3 & Sound Tracker v3.x Compiled player & Player decompiler \\
COMPILEDSTP1 & Sound Tracker Pro v1.x player & Player decompiler \\
COMPILEDSTP2 & Sound Tracker Pro v2.x player & Player decompiler \\
RAW & Raw scaner & Raw data scanner \\
SNA128 & SNA 128k & Snapshot \\
Z80V145 & Z80 v1.45 & Snapshot \\
Z80V20 & Z80 v2.x & Snapshot \\
Z80V30 & Z80 v3.x & Snapshot \\
ZXSTATE & SZX (ZX-State) & Snapshot \\
\bottomrule \end{longtable}}


\begin{document}

\frontmatter
\begin{titlepage}
\pagecolor{lhunanight}\color{white}
\centering
\vspace*{3cm}
\includegraphics[width=0.6\textwidth]{lhuna_logo}\par
\vspace{1.5cm}
{\Huge\bfseries ZXTune\_LHUNA\par}\vspace{0.4cm}{\LARGE Manuel de référence\par}
\vspace{0.5cm}
{\Large Lecteur de musiques de puces sonores pour Android\par}
\vfill
{\large Version du document : \today\par}
\vspace{1cm}
\end{titlepage}
\nopagecolor\color{black}

\chapter*{Licence et marque}
\addcontentsline{toc}{chapter}{Licence et marque}

\section*{Ce qui est sous licence LHUNA}
Le nom \emph{LHUNA}, le logo, l'identité visuelle et le présent manuel sont la propriété de leur auteur et sont
diffusés sous licence LHUNA. Toute reproduction du logo ou de ce manuel demande son accord.

\section*{Ce qui est sous licence libre}
L'application \soft{} est un logiciel libre distribué sous \textbf{GNU General Public License, version 3}. Ce choix
n'est pas arbitraire : il découle des composants qu'elle contient.
\begin{itemize}
  \item \soft{} est dérivé de \emph{ZXTune}, distribué sous GNU Lesser General Public License version~3.
        La LGPL~v3 n'autorise une combinaison qu'avec la GPL~v3, pas avec la GPL~v2.
  \item Plusieurs moteurs d'émulation intégrés sont sous GPL~v3 (cœurs Saturn et PlayStation
        \emph{Highly Theoretical/Experimental}) ou GPL~v2 « ou ultérieure » (libsidplayfp, z80ex, ASAP), compatibles
        avec la GPL~v3.
  \item Les autres composants sont sous licences permissives (MIT, BSD, ISC, zlib, MPL~2.0…), toutes compatibles.
\end{itemize}
Conséquence : quiconque reçoit l'APK a le droit d'obtenir le code source correspondant, de le modifier et de le
redistribuer sous les mêmes conditions. La liste complète des composants et de leurs licences figure en
annexe~\ref{chap:credits}.

\section*{Les drivers propriétaires}
Les drivers musicaux de l'auteur (famille \emph{NebulaWeave}) ne font pas partie de l'APK de base. Ils sont
distribués dans un APK séparé, sous licence propriétaire.
\begin{itemize}
  \item Sans cet APK, \soft{} fonctionne exactement comme décrit dans ce manuel.
  \item Avec cet APK installé, \soft{} peut utiliser les drivers qu'il fournit.
\end{itemize}

\begin{technique}
Pour que les deux licences cohabitent, l'APK des drivers doit rester une œuvre séparée au sens de la GPL~v3 :
il fonctionne dans son propre processus et échange avec \soft{} par une interface documentée (service Android,
fournisseur de contenu, données). Charger sa bibliothèque native dans le processus de \soft{} ferait de l'ensemble une
seule œuvre, qui devrait alors être entièrement sous GPL~v3.
\end{technique}

\section*{Marques citées}
Les noms de consoles, d'ordinateurs, de puces et de jeux cités dans ce manuel sont la propriété de leurs détenteurs
respectifs et ne sont utilisés que pour décrire la compatibilité.

\tableofcontents
\chapter*{Comment lire ce manuel}
\addcontentsline{toc}{chapter}{Comment lire ce manuel}
Ce manuel s'adresse à deux publics à la fois : l'auditeur qui veut simplement écouter ses musiques, et le
musicien ou passionné de matériel qui veut comprendre exactement ce que fait chaque puce.

\begin{enclair}
Les encadrés comme celui-ci expliquent une notion sans vocabulaire technique. On peut lire tout le manuel en ne
suivant que le texte courant et ces encadrés.
\end{enclair}

\begin{technique}
Les encadrés comme celui-ci donnent le détail matériel : registres, fréquences, formules. On peut les sauter sans
perdre le fil.
\end{technique}

Les captures d'écran ont été produites par l'application elle-même, sur un téléphone de 411$\times$915~dp. Sur un
autre appareil, les proportions peuvent varier, mais les éléments restent les mêmes.


\mainmatter
\part{Découvrir LHUNA}
\chapter{Présentation}

\section{Ce qu'est \soft{}}
\soft{} est un lecteur de musique pour Android consacré aux musiques des ordinateurs, consoles et bornes d'arcade
des années 1980 à 2000. Ces musiques ne sont pas des enregistrements : ce sont des programmes et des données qui
pilotent une puce sonore. Pour les écouter fidèlement, \soft{} reproduit, par logiciel, le comportement de ces puces.
On parle d'\emph{émulation}.

\begin{enclair}
Un fichier MP3 contient le son lui-même. Un fichier de musique de Commodore~64 ou de Mega Drive contient plutôt une
partition et les réglages d'un instrument électronique précis. \soft{} joue cette partition sur une copie logicielle
de l'instrument d'origine.
\end{enclair}

\section{Origine}
\soft{} est dérivé de \emph{ZXTune}, lecteur libre créé pour les musiques du ZX Spectrum puis étendu à de nombreuses
plateformes. \soft{} en reprend le cœur de lecture et y ajoute :
\begin{itemize}
  \item des moteurs d'émulation choisis pour leur précision (par exemple ares pour la Super Nintendo, Ootake pour la
        PC Engine, reSIDfp pour le Commodore~64) ;
  \item une visualisation en temps réel de chaque voix et de l'état interne des puces, puce par puce ;
  \item une interface repensée : listes de lecture multiples, favoris, file d'attente, réglages en direct.
\end{itemize}

\section{Ce que \soft{} sait faire}
\begin{itemize}
  \item Lire \lhunaModulesCount{} formats de musique et ouvrir \lhunaContainersCount{} formats de conteneurs
        (archives, images de disquettes, compresseurs), voir la partie~\ref{part:compat}.
  \item Parcourir le stockage du téléphone et de grands catalogues en ligne.
  \item Afficher un oscilloscope par voix, la réponse du filtre du SID, et des tableaux de bord détaillés pour les
        puces SID, FM, SNES, PC Engine, NES, MT-32 et les puces PCM des bornes Sega.
  \item Fonctionner avec Android Auto et l'écran verrouillé.
\end{itemize}

\section{Configuration requise}
\begin{itemize}
  \item Android~4.1 ou plus récent pour la version complète.
  \item Pour l'oscilloscope : OpenGL ES, ou Vulkan quand l'appareil le permet (plus fluide).
  \item Espace de stockage : environ 16~Mo pour l'application, plus vos musiques.
\end{itemize}

\section{Historique}
\subsection{ZXTune}
ZXTune est né dans la communauté du ZX Spectrum, dont la scène musicale a produit des dizaines de milliers de
morceaux pour la puce AY-3-8910 avec des trackers comme Sound Tracker, Pro Tracker ou ASC Sound Master. Chacun de
ces logiciels enregistrait ses musiques dans son propre format, et les lecteurs de l'époque ne savaient en lire
qu'un seul. ZXTune, créé par Vitamin, s'est donné pour but de tous les lire, de les retrouver même cachés dans des
images de disquettes, des archives ou des programmes compressés, et de les rejouer fidèlement.

Le projet s'est ensuite ouvert aux autres plateformes en intégrant les meilleures bibliothèques libres existantes :
libsidplayfp pour le Commodore~64, Game Music Emu pour les consoles, libopenmpt et libxmp pour les modules de
tracker, vgmstream pour l'audio des jeux récents, libvgm pour les enregistrements VGM. Il existe en ligne de
commande, avec une interface graphique pour ordinateur et en application Android.

\subsection{ZXTune\_LHUNA}
\soft est une branche de l'application Android de ZXTune, centrée sur deux objectifs :
\begin{description}
  \item[La fidélité] choisir pour chaque puce le moteur d'émulation le plus juste disponible, et pouvoir comparer
        avec un moteur plus léger quand l'appareil peine ;
  \item[La lisibilité] montrer ce que fait chaque puce, voix par voix et image par image, avec des vues pensées
        pour la synthèse propre à chaque puce plutôt qu'une vue générique.
\end{description}
Tout ce qui est affiché est l'état réel de la puce émulée à l'instant présent : pas de valeurs moyennées ni
inventées. Le texte complète toujours le graphique, pour qu'un auditeur non technicien comprenne ce qu'il voit et
qu'un musicien retrouve exactement les valeurs de ses registres.

\chapter{Installation et prise en main}

\section{Installer l'application}
\soft est distribué sous forme de fichier APK. Copiez-le sur votre téléphone ou téléchargez-le, puis ouvrez-le.
Android vous demande d'autoriser l'installation depuis cette source : acceptez pour cette fois. Les mises à jour
s'installent de la même façon, par-dessus la version précédente ; vos listes et vos réglages sont conservés.

\section{Les autorisations}
Au premier lancement, \soft vous demande l'accès à vos musiques. Sur les versions récentes d'Android, vous
choisissez un dossier précis plutôt que de donner accès à tout le téléphone ; nous vous conseillons de désigner le
dossier qui contient votre collection. L'autorisation des notifications sert à afficher la lecture en cours et la
progression quand vous ajoutez un dossier entier à une liste.

\section{Les trois écrans}
L'application tient en trois écrans, accessibles par la barre du bas. L'écran \textbf{Lecture} montre le morceau en
cours, l'oscilloscope, le tableau de bord de la puce et les commandes. L'écran \textbf{Listes} rassemble vos listes
de lecture. L'écran \textbf{Navigateur} donne accès à vos fichiers et aux catalogues en ligne.

\capturedouble{now-playing-en}{playlist-en}{Écran Lecture}{écran Listes}

Pour votre première écoute, ouvrez le navigateur, choisissez un catalogue (la High Voltage SID Collection est un bon
point de départ) et touchez un morceau. La lecture démarre et l'écran Lecture se remplit : en haut l'oscilloscope
de chaque voix, en dessous le tableau de bord de la puce. Si vous préférez l'écran épuré, repliez les informations
avec la flèche en bas à gauche.

\capture{browser-sources-en}{Navigateur : les sources}

\chapter{L'interface}

\section{Écran Lecture}
\capture{now-playing-en}{Écran Lecture, informations repliées}

De haut en bas :
\begin{description}
  \item[Barre d'outils] logo, plein écran, réglages en direct, choix de la visualisation (oscilloscope ou réponse
        du filtre), tableau de bord, menu.
  \item[Visualisation] oscilloscope par voix et, en dessous, le tableau de bord de la puce (chapitre~\ref{chap:dashboards}).
  \item[Informations] titre, auteur, puce, format et pilote d'origine. Repliées par défaut pour laisser la place à
        la visualisation ; la flèche en bas à gauche les déplie.
  \item[Position] temps écoulé, barre de progression, durée totale.
  \item[Commandes] informations, aléatoire, précédent, lecture/pause, suivant, répétition, bouton \textbf{$\equiv$+}
        (favoris, ajouter à une liste), menu du morceau (partager, envoyer, sonnerie, propriétés).
\end{description}

\capture{now-playing-info-en}{Écran Lecture, informations dépliées}

\section{Écran Listes}
\capture{playlist-en}{Écran Listes}
En haut, la recherche et les actions sur la liste (statistiques, export, effacer, trier). Sous la recherche, le
sélecteur de liste indique la liste affichée et son nombre de morceaux ; le toucher ouvre le menu des listes
(chapitre~\ref{chap:listes}). Chaque morceau montre la famille de puce, le titre, l'auteur et la durée ; la poignée
à droite permet de le déplacer.

\section{Écran Navigateur}
\capturedouble{browser-sources-en}{browser-folder-en}{Sources}{contenu d'un dossier}
Le chemin parcouru s'affiche en haut ; chaque élément permet de revenir directement à ce niveau.

\section{Mini-lecteur et barre de navigation}
Hors de l'écran Lecture, un mini-lecteur rappelle le morceau en cours et permet de passer au suivant ou de mettre
en pause.

\section{Paysage et tablette}
\capture[0.9]{landscape-now-playing-en}{Écran Lecture en paysage}
\capture[0.9]{tablet-playlist-en}{Sur tablette}

\section{Plein écran}
Le bouton plein écran affiche l'oscilloscope seul, et en portrait le tableau de bord sous lui. En portrait,
l'affichage évite l'encoche de la caméra en haut et les coins arrondis en bas ; en paysage, l'encoche latérale.
Un appui court quitte le plein écran, un appui long ouvre les réglages en direct.


\part{Écouter}
\chapter{Le navigateur}

\section{Sources locales}
Le navigateur donne accès au stockage du téléphone et aux cartes SD. Sur Android~10 et plus, l'accès passe par le
sélecteur de dossiers du système : vous choisissez une fois le dossier de vos musiques, LHUNA s'en souvient.

\section{Catalogues en ligne}
LHUNA donne accès à plusieurs catalogues publics de musiques de puces, consultés en direct :
\begin{itemize}
  \item High Voltage SID Collection (Commodore~64) ;
  \item The Mod Archive, Modland et Amiga Music Preservation (modules de tracker, Amiga) ;
  \item ZXArt et ZXTunes (ZX Spectrum), Aygor ;
  \item Atari SAP Music Archive (Atari 8 bits) ;
  \item VGMRips (enregistrements VGM de consoles et bornes d'arcade) ;
  \item Joshw (rips de consoles), Khinsider, OverClocked ReMix, Aminet, Scene.org ;
  \item des radios en ligne.
\end{itemize}
\arediger{Description de chaque catalogue : contenu, taille, mainteneurs, conditions d'utilisation. À vérifier sur
leurs sites.}

\section{Archives et fichiers multi-morceaux}
LHUNA ouvre les archives (ZIP, RAR, 7z, LHA\dots) et les images de disquettes (TR-DOS, SCL\dots) comme des dossiers.
Un fichier qui contient plusieurs morceaux (SID, NSF, AY, KSS\dots) est en revanche lu directement : ses
sous-morceaux s'enchaînent, puis la lecture continue avec le fichier suivant du dossier.

\section{Recherche et sélection}
La loupe filtre le dossier affiché. Un appui long sur un élément démarre la sélection multiple ; la barre du haut
propose alors :
\begin{description}
  \item[Ajouter] demande la liste de destination et y ajoute les fichiers (les dossiers sont parcourus en entier) ;
  \item[Jouer ensuite] place les éléments sélectionnés dans la file d'attente (section~\ref{sec:file}).
\end{description}

\chapter{Listes de lecture}\label{chap:listes}

\section{Plusieurs listes}
Vous pouvez créer autant de listes que vous voulez. Deux existent d'emblée : la liste par défaut, qui reprend votre
ancienne liste lors de la mise à jour, et les favoris, que vous remplissez avec le bouton $\equiv$+ de l'écran
Lecture.

La liste affichée est aussi celle qui se lit. Quand vous changez de liste avec le sélecteur, nous enregistrons la
liste que vous quittez et chargeons la nouvelle ; le morceau en cours continue, et le bouton « suivant » suivra la
nouvelle liste. Nous avons fait ce choix pour que l'écran montre toujours ce que vous allez entendre.

\capture{playlist-en-top}{Le sélecteur de liste}

\section{Gérer vos listes}
Le menu du sélecteur permet de créer une liste, de renommer ou dupliquer la liste affichée, de la supprimer (une
autre liste prend alors sa place) et d'échanger des listes au format M3U.

\section{Ajouter des morceaux}
Depuis l'écran Lecture, le bouton $\equiv$+ puis « Add to playlist\dots » range le morceau en cours ; vous pouvez
aussi créer une liste au passage. Depuis le navigateur, sélectionnez des fichiers ou des dossiers puis touchez
« Add ». Depuis une liste, sélectionnez des morceaux et copiez-les dans une autre liste.

\section{La file d'attente}\label{sec:file}
« Play next » insère des morceaux juste après celui en cours, sans toucher à vos listes. Ils passent dans l'ordre où
vous les avez ajoutés, puis la lecture reprend là où elle en était. Si vous choisissez vous-même un autre morceau
entre-temps, la file est conservée mais ce point de reprise est oublié.

\section{Import et export M3U}
L'export écrit la liste affichée au format M3U étendu, avec titre, auteur et durée. Les emplacements sont complets,
y compris à l'intérieur des archives et des catalogues en ligne : une liste exportée se réimporte à l'identique.
L'import accepte aussi les listes d'autres lecteurs et crée une nouvelle liste portant le nom du fichier. Vos
anciennes listes au format XSPF restent accessibles dans le navigateur.

\section{Trier et ranger}
Le tri porte sur le titre, l'auteur ou la durée. Pour un ordre à votre goût, déplacez les morceaux par leur poignée.
Les statistiques donnent le nombre de morceaux, le nombre de fichiers distincts et la durée totale.

\chapter{La lecture}

\section{Aléatoire et répétition}
La lecture aléatoire choisit le morceau suivant au hasard dans la liste ou le dossier. La répétition rejoue le
morceau en cours sans fin ; sans elle, le morceau s'arrête à la fin de sa durée connue.

\section{Durée, boucles et fondus}
Beaucoup de musiques de jeux bouclent indéfiniment : le fichier ne dit pas quand s'arrêter. Quand une durée est
connue, par le format lui-même ou par une base de données (la base \emph{Songlengths} de la HVSC pour le SID), nous
l'utilisons. Sinon, nous appliquons la durée par défaut, que vous réglez dans les réglages, ainsi que la limite de
boucles et les fondus. Quand un morceau a une introduction avant sa boucle, la barre de progression marque le début
de la partie bouclée.

\section{Hors de l'application}
Vous pouvez piloter la lecture depuis la notification, l'écran verrouillé, un casque Bluetooth ou Android Auto.
Android Auto affiche vos listes et une visualisation simplifiée.

\section{Partager}
Le menu du morceau partage le fichier, envoie un lien vers le morceau, crée une sonnerie à partir d'un extrait et
affiche toutes les propriétés connues du morceau.


\part{Compatibilité}\label{part:compat}
\chapter{Plateformes et puces émulées}

\begin{enclair}
Chaque machine de l'époque avait sa propre puce sonore, avec ses forces et ses limites : trois voix carrées pour le
ZX Spectrum, un synthétiseur analogique pour le Commodore~64, des opérateurs FM pour la Mega Drive, des échantillons
pour la Super Nintendo. Ce chapitre présente chaque famille et ce que LHUNA en affiche.
\end{enclair}

\begin{longtable}{@{}p{0.22\textwidth} p{0.27\textwidth} p{0.43\textwidth}@{}}
\toprule \textbf{Plateforme} & \textbf{Puce} & \textbf{Synthèse} \\ \midrule \endhead
ZX Spectrum 128, Amstrad CPC, Atari ST, MSX & AY-3-8910 / YM2149 & 3 voix carrées, bruit, enveloppe matérielle \\
ZX Spectrum (TurboSound) & 2 $\times$ AY & 6 voix \\
ZX Spectrum (TurboFM) & 2 $\times$ YM2203 & FM 4 opérateurs + AY \\
SAM Coupé & SAA1099 & 6 voix carrées, bruit, enveloppes \\
Commodore 64 & MOS 6581 / 8580 (SID), jusqu'à 3 puces & 3 oscillateurs, filtre analogique, modulation en anneau, synchronisation \\
NES / Famicom & Ricoh 2A03 / 2A07 & 2 carrés, triangle, bruit, DPCM \\
Super Nintendo & SPC700 + S-DSP & 8 voix d'échantillons BRR, ADSR, écho FIR \\
Game Boy & LR35902 (APU) & 2 carrés, onde programmable, bruit \\
PC Engine / TurboGrafx-16 & HuC6280 & 6 voix à table d'ondes 32 $\times$ 5 bits, bruit, DDA, LFO \\
Atari 8 bits & POKEY & 4 voix, polynômes de bruit, filtres \\
Mega Drive / Genesis & YM2612 + SN76489 & FM 6 voix 4 opérateurs, DAC, PSG \\
PC (AdLib, Sound Blaster) & YM3526 / YM3812 / YMF262 & FM 2 ou 4 opérateurs, mode rythme \\
Sharp X68000, bornes & YM2151 & FM 8 voix 4 opérateurs \\
Bornes Sega Model 1 et 2 & YMW258 (MultiPCM) & 28 voix d'échantillons \\
Sega Saturn, bornes Model 2A/3 & YMF292 (SCSP) & 32 voix PCM/FM, DSP d'effets \\
Roland MT-32 & LA (Munt) & synthèse linéaire arithmétique \\
\bottomrule
\end{longtable}

\arediger{Une section par plateforme : historique de la puce, concepteurs, années, musiciens et jeux marquants,
particularités audibles, avec une capture du tableau de bord correspondant. Recherche de sources en cours.}

\chapter{Formats de musique}
\soft{} reconnaît \lhunaModulesCount{} formats de musique. Nous avons généré la table ci-dessous à partir de
l'application elle-même ; elle est donc exacte pour la version décrite.

Derrière ce nombre se cachent quatre façons très différentes de conserver une musique. Savoir à laquelle appartient
un fichier vous dit ce que vous pouvez en attendre : un module de tracker se rejoue toujours à l'identique, un
enregistrement de registres ne contient que ce qui a été capté, un rip exécute le vrai programme de la machine et
réagit donc à l'émulation dans ses moindres détails.

\begin{description}
  \item[Tracker module] partition et instruments, rejoués par un lecteur fidèle à celui d'origine ;
  \item[Register stream] suite de valeurs écrites dans les registres de la puce, enregistrées depuis la machine ;
  \item[Memory dump] programme de lecture d'origine exécuté par un processeur émulé (Z80, 6502\dots) ;
  \item[Multichip] plusieurs puces jouées ensemble.
\end{description}


{\small\lhunaFormatsTable}

\section{Conteneurs}
\soft{} ouvre aussi \lhunaContainersCount{} formats de conteneurs : archives, compresseurs, images disque,
instantanés de machine et décompilateurs de lecteurs.

{\small\lhunaContainersTable}

\chapter{Formats audio enregistrés}
En plus des musiques émulées, LHUNA lit les formats audio courants :
\begin{itemize}
  \item WAV : PCM 8 à 32 bits, virgule flottante 32 bits, en-tête étendu (WAVE\_FORMAT\_EXTENSIBLE) ;
  \item FLAC, Monkey's Audio, Musepack, TAK (sans perte) ;
  \item MP3, Ogg Vorbis, Opus, AAC, AC-3 ;
  \item et les formats audio de jeux pris en charge par vgmstream (chapitre précédent).
\end{itemize}
Pour ces formats, l'écran Lecture affiche le codec, la fréquence d'échantillonnage, la résolution et le débit.
\capture{now-playing-info-flac-fr}{Informations d'un fichier FLAC}


\part{Moteurs d'émulation}
\chapter{Moteurs d'émulation}

\begin{enclair}
Plusieurs programmes savent imiter une même puce. Certains sont très précis mais demandent beaucoup de calcul ;
d'autres sont plus légers. Quand il existe un vrai choix, LHUNA le propose dans les réglages.
\end{enclair}

\section{Commodore 64 : reSID et reSIDfp}
Les deux moteurs proviennent de la bibliothèque libsidplayfp, qui émule aussi le processeur et les composants du
Commodore~64 nécessaires à l'exécution des lecteurs d'origine.
\begin{description}
  \item[reSIDfp] (par défaut) modèle à virgule flottante du filtre analogique et des formes d'onde combinées ;
  \item[reSID] modèle historique, plus léger.
\end{description}
Le modèle de puce (6581 ou 8580), les courbes de filtre et le biais se règlent (chapitre~\ref{chap:emu-settings}).

\section{Super Nintendo : ares ou snes\_spc}
\begin{description}
  \item[ares] (par défaut) émulation au cycle près du processeur SPC700 et du S-DSP, issue de l'émulateur ares ;
        la stéréo est rendue comme sur la console, inversions de phase comprises ;
  \item[snes\_spc] moteur historique de Shay Green (blargg), plus léger.
\end{description}
Le même réglage s'applique aux fichiers SPC et aux rips SNSF, qui exécutent aussi le processeur principal 65816.

\section{Puces FM}
\begin{longtable}{@{}l l p{0.5\textwidth}@{}}
\toprule \textbf{Famille} & \textbf{Puces} & \textbf{Cœurs proposés} \\ \midrule
OPN2 & YM2612, YM3438 & Nuked (précis), ymfm, Genesis Plus GX (rapide) \\
OPM & YM2151 & Nuked (précis), ymfm, MAME \\
OPN & YM2203, YM2608, YM2610 & ymfm, MAME \\
OPL & YM3526, YM3812, YMF262 & Nuked (précis), AdLibEmu, MAME \\
OPLL & YM2413 & Nuked (précis), EMU2413, MAME \\
\bottomrule
\end{longtable}

\section{PC Engine : Ootake}
La puce HuC6280 est émulée par le cœur d'Ootake, émulateur PC Engine japonais réputé pour sa fidélité. LHUNA
l'utilise pour les enregistrements VGM et y lit la table d'ondes de chaque voix pour le tableau de bord.

\section{Sega Saturn et bornes : SCSP}
Le SCSP est émulé avec son DSP d'effets activé : la réverbération et les effets programmés par le jeu sont audibles.

\arediger{Historique et auteurs de chaque moteur (Dag Lem, Antti Lankila et Leandro Nini pour reSID/reSIDfp ;
Near et l'équipe ares ; nukeykt pour les cœurs Nuked ; Aaron Giles pour ymfm ; Kitao Nakamura pour Ootake\dots),
mesures de précision, coût en calcul sur téléphone.}

\chapter{Réglages d'émulation}\label{chap:emu-settings}

\section{SID}
\capture{settings-sid}{Réglages SID dans le panneau en direct}
\begin{description}
  \item[Moteur] reSID ou reSIDfp.
  \item[Modèle] 6581 ou 8580 pour les morceaux qui ne le précisent pas, ou pour tous si le choix est forcé.
  \item[Filtre] préréglages mesurés sur de vraies puces, ou courbe personnalisée.
  \item[Sortie 2SID/3SID] stéréo (puces réparties à gauche, à droite, au centre) ou mono.
  \item[Voie des samples (\$D418)] ajoute à l'oscilloscope une voie par puce montrant les échantillons joués par le
        registre de volume ; les trois voix ne les contiennent alors plus.
\end{description}

\section{SNES}
\capture{settings-spc}{Réglages SNES}
\begin{description}
  \item[Moteur d'émulation] ares ou snes\_spc, appliqué au morceau suivant.
\end{description}

\section{FM}
\capturedouble{settings-ym2612-psg}{settings-ym3812}{Réglages OPN2}{réglages OPL}
Pour chaque famille : cœur d'émulation, préréglage de puce, échelle (« ladder ») du DAC, filtres passe-bas et
passe-haut de sortie, largeur stéréo.

\section{AY et Z80}
\capture{settings-ay}{Réglages AY}
\begin{description}
  \item[Fréquence d'horloge] doit correspondre à la machine d'origine : une horloge fausse décale toutes les notes.
  \item[Type de puce] AY-3-8910 ou YM2149 : le YM a une enveloppe plus fine et des volumes légèrement différents.
  \item[Disposition stéréo] répartition des trois voix (ABC, ACB, mono\dots), selon les habitudes de chaque
        plateforme.
  \item[Interpolation] adoucit le rendu ; désactivée, le son est plus brut, proche de la puce.
\end{description}

\section{Effet des réglages FM}
\begin{description}
  \item[Ladder] les YM2612 d'origine ont un défaut de leur convertisseur qui renforce les sons faibles ; certaines
        bandes-son ont été composées en en tenant compte. Le YM3438 et les Mega Drive tardives ne l'ont pas.
  \item[Filtres de sortie] reproduisent le filtrage analogique de la console ; sans eux, le son est plus brillant
        que sur le matériel.
  \item[Stéréo] largeur de la stéréo : à 0, mono.
\end{description}

\section{Durées et volume}
\begin{description}
  \item[Durée par défaut] utilisée quand ni le format ni une base de données ne donnent la durée.
  \item[Détection du silence] passe au morceau suivant après un silence de la durée choisie.
  \item[Limite de boucles] nombre de répétitions des morceaux qui bouclent.
  \item[Fondus] fondu d'entrée et de sortie.
  \item[Gain] volume général, de 50 à 200~\%.
\end{description}
La liste complète, avec les valeurs par défaut, figure au chapitre~\ref{chap:all-settings}.


\part{Visualiser}
\chapter{L'oscilloscope}

\section{Une voie par voix}
L'oscilloscope dessine la forme du son de chaque voix. Chaque puce occupe une colonne, chaque voix une ligne, avec
son nom en haut de la case. Vous voyez ainsi qui joue quoi : la basse qui tient, l'arpège qui scintille, la caisse
claire qui passe en un éclair.

Pour que les formes restent immobiles, nous calons la fenêtre de chaque voie sur la période de sa note : c'est le
déclenchement, comme sur un oscilloscope de laboratoire. Nous réduisons ensuite le signal à deux points par colonne
de pixels, le minimum et le maximum, si bien qu'aucun pic n'est perdu même quand la voie est étroite.

\section{Cas particuliers}
Avec la voie des samples du SID activée, une quatrième ligne par puce montre les échantillons joués par le volume.
Pour l'AdLib, nous disposons les neuf voix mélodiques en 3 sur 3 et masquons les cinq voies du mode rythme, que peu
de morceaux utilisent ; activez « AdLib rhythm lanes » si vous en avez besoin.

\section{La réponse du filtre du SID}
Le bouton de visualisation remplace la moitié haute de l'oscilloscope par la courbe du filtre de chaque SID, avec la
position de chaque voix. Les états précédents restent visibles quelques instants en fondu, ce qui rend lisibles les
balayages rapides de la coupure.
\capturedouble{ux4-lowpass-waves}{ux4-sweep}{Filtre passe-bas résonant}{balayage de la coupure}
Si vous voulez plus de place pour les voix, réduisez la part de la réponse du filtre dans les réglages de
l'oscilloscope.
\capture[0.7]{ux4-lowpass-share35}{Réponse du filtre réduite à 35~\% de la hauteur}

\section{Réglages}
\capture{prefs-en-oscilloscope}{Réglages de l'oscilloscope}
Vous réglez la fenêtre de temps, le gain, le nombre de périodes affichées, la hauteur de l'oscilloscope, la part de la
réponse du filtre, la fréquence d'images maximale et le moteur de rendu. Vulkan donne le rendu le plus fluide quand
votre téléphone le permet ; sinon, OpenGL ES prend le relais.

\chapter{Tableaux de bord par puce}\label{chap:dashboards}

L'oscilloscope vous montre la forme du son ; le tableau de bord, en dessous, vous dit comment la puce la fabrique.
Nous n'avons pas voulu une vue unique remplie de registres : chaque famille de puces a la sienne, construite autour de
ce qui fait son son. Sur le SID, c'est le filtre et les modulations entre voix ; sur la Super Nintendo, les
enveloppes et l'écho ; sur la PC Engine, la table d'ondes ; sur les puces FM, les opérateurs et la façon dont ils
sont branchés.

Nous relevons l'état de la puce environ toutes les millisecondes de son émulé, et chaque image affiche le dernier
relevé, tel quel. Si une valeur saute, c'est que la puce a sauté.

\section{SID (Commodore 64)}
\capture[0.8]{sidcells-now-playing}{Tableau de bord SID}
Chaque colonne correspond à une voix. De haut en bas, vous lisez :
\begin{description}
  \item[Wave] la forme d'onde (triangle, dent de scie, impulsion, bruit, ou une combinaison) et, pour l'impulsion,
        son rapport cyclique. Quand le rapport cyclique bouge sans arrêt, c'est la « PWM » qui donne leur
        épaisseur aux basses de Rob Hubbard.
  \item[Env] l'enveloppe : le niveau actuel, la phase (attaque, déclin, maintien, relâchement) et les durées ADSR en
        millisecondes. Une attaque à 2~ms et un déclin court, c'est une note percussive ; un maintien haut, une note
        tenue.
  \item[Freq] la note la plus proche et la fréquence en hertz. Un arpège se voit ici : la note change à chaque image.
  \item[Mod] la synchronisation et la modulation en anneau. Nous n'indiquons que la voix source : « Sync V3 »,
        « Ring V2 », « Sync+Ring V1 ». Sur le SID, une voix est toujours modulée par la précédente (V1 par V3, V2 par
        V1, V3 par V2) ; connaître la source suffit.
\end{description}
La dernière ligne montre le volume général, la résonance, le type de filtre et les voix qui passent dedans. Si vous
voyez une voix entrer et sortir du filtre au rythme de la musique, vous tenez l'un des trucs favoris des musiciens
de l'époque.
\capturedouble{sidcells-dual}{sidcells-fullscreen}{Deux puces SID}{plein écran}

\section{Super Nintendo (S-DSP)}
\capturedouble{snes}{snes-inspector}{Vue des 8 voix}{inspecteur d'une voix}
Le S-DSP joue huit voix d'échantillons. Pour chacune, vous voyez le numéro de l'échantillon, la note, l'enveloppe
dessinée (ADSR ou mode GAIN, que beaucoup de drivers utilisent pour leurs fondus maison), les volumes gauche et
droite, et trois indicateurs : écho, bruit, modulation de hauteur.

En bas, les volumes principaux, le volume et le retour de l'écho, et les huit coefficients du filtre FIR qui colore
ses répétitions. Un retour élevé avec un FIR passe-bas, c'est la grande salle réverbérante des jeux Square ; un écho
court et sec, une ambiance de donjon. Touchez une voix pour ouvrir son inspecteur et voir son enveloppe en grand.

\section{PC Engine (HuC6280)}
\capture[0.8]{pce-dashboard}{Vue PC Engine}
Chacune des six voix lit en boucle une table de 32 valeurs de 5 bits : sa forme d'onde. Le rectangle de chaque voix
dessine la table en cours d'utilisation, pas à pas. À côté, une barre de niveau et deux barres de panoramique, puis
la note, la fréquence et le niveau.

Nous décrivons la forme en mots : « Sine », « Triangle », « Saw $\uparrow$ », « Square », « Pulse 25\% »\dots{} et
nous numérotons chaque table distincte. Quand le numéro change en plein morceau, le driver vient de recharger la
voix avec un autre instrument. Le HuC6280 ne sait pas lire une table à l'envers : une scie descendante est une table
que le driver a écrite à l'envers, d'où la mention « Saw $\downarrow$ (reversed) ».

La ligne du bas de chaque voix donne ses registres tels qu'ils ont été écrits : \$802 et \$803 (fréquence),
\$804 (volume et activation), \$805 (panoramique), \$806 (mode DDA), \$807 (bruit). La ligne globale donne \$801
(volume principal), \$808 et \$809 (LFO). Le mode DDA, quand il est actif, montre que la voix joue des échantillons
envoyés un par un par le processeur : c'est ainsi que la PC Engine fait parler ses personnages.

\section{Puces FM}
\capturedouble{fm-megadrive-cards}{fm-megadrive-inspector}{Mega Drive : cartes des voix}{inspecteur d'une voix}
Chaque carte montre une voix : son algorithme, dessiné avec les porteuses en bas et les modulateurs au-dessus, et le
niveau de chaque opérateur. Retenez une règle simple : plus un modulateur est fort, plus le son est riche et dur.
Une basse de Mega Drive qui « claque » a un modulateur dont l'enveloppe chute très vite.

Touchez une voix pour ouvrir l'inspecteur : enveloppes de chaque opérateur dessinées en grand, multiplicateurs de
fréquence, désaccord, rétroaction de l'opérateur 1 et sensibilité au LFO.
\capturedouble{fmview-algorithms}{fm-adlib-cards}{Algorithmes FM}{AdLib}
Pour l'AdLib (OPL2), chaque voix n'a que deux opérateurs et deux branchements possibles : en série (l'un module
l'autre) ou en parallèle (les deux sont entendus). Nous disposons les neuf voix en 3 sur 3.

\section{NES, MT-32 et puces PCM}
\capturedouble{nes}{mt32parts}{NES}{Roland MT-32 : parties}
Sur NES, les cartes des deux voix carrées dessinent leur rapport cyclique (12,5, 25, 50 ou 75~\%) et donnent le
volume sur 4 bits, l'enveloppe et le balayage de hauteur (« sweep ») ; le triangle et le bruit suivent, puis l'état
du canal DMC, qui joue les échantillons. Pour le MT-32, il n'y a pas de registres utiles à montrer : chaque partie
affiche le nom de son instrument, les touches jouées et l'état de ses « partiels », les briques de synthèse que le
module combine pour former un son.
\capturedouble{pcmchips-model1-portrait}{pcmchips-scsp-portrait}{YMW258 (Model 1/2)}{SCSP (Saturn, Model 2A/3)}
Les puces PCM des bornes Sega et de la Saturn jouent des dizaines d'échantillons à la fois. Chaque case montre une
voix active : son mode (PCM, FM ou bruit sur le SCSP), sa note, son panoramique et son niveau total (TL) en décibels.

\chapter{Lire les valeurs}

Quelques conventions valent pour toutes les vues ; une fois que vous les connaissez, vous lisez n'importe quel
tableau de bord.

\section{Ce que vous voyez est l'état exact}
Nous n'affichons aucune moyenne et aucun lissage. Une valeur moyennée est une valeur que la puce n'a jamais eue, et
elle cache justement ce qui nous intéresse : un arpège qui change de note à chaque image, une enveloppe qui
redémarre. Si une valeur scintille, c'est que la musique la fait scintiller.

\section{Notes et fréquences}
Les notes sont écrites à l'anglaise : C, D, E, F, G, A, B, avec \# pour les dièses. L'octave 4 contient le la à
440~Hz (A4), le do du milieu du piano est C4. Quand la fréquence tombe entre deux notes, nous affichons la plus
proche ; la fréquence en hertz vous donne l'écart. Un léger désaccord entre deux voix qui jouent la même note est
souvent voulu : c'est ce qui épaissit le son.

\section{Niveaux}
Les niveaux sont exprimés par rapport au maximum de la puce : 0~dB est le plein volume, $-6$~dB la moitié de
l'amplitude, $-\infty$ le silence. Les barres suivent la même échelle.

\section{Registres}
Les registres sont en hexadécimal, précédés de \$, comme dans les documentations d'époque : \$D418 sur le SID,
\$804 sur la PC Engine. Nous les montrons tels qu'ils ont été écrits par le driver, pas tels que nous les
interprétons ; si vous cherchez dans une documentation, vous retrouverez les mêmes valeurs.

\chapter{Panneau de réglages en direct}
Changer un réglage d'émulation, puis revenir au morceau pour l'écouter, ne permet pas de comparer : l'oreille oublie
vite. Le panneau en direct règle ce problème. Un appui long sur l'oscilloscope, ou le bouton des réglages, l'ouvre
par-dessus la visualisation pendant que la musique continue.

Il ne montre que ce qui concerne le morceau en cours : les réglages de la puce jouée, ceux de l'oscilloscope et les
réglages généraux. Vous pouvez donc passer de reSIDfp à reSID, ou d'un filtre 6581 à un autre, et entendre la
différence dans la seconde. Les rares réglages qui demandent de recharger le morceau sont signalés par la mention
« applied to the next track ».
\capturedouble{settings-sid}{settings-general}{Page SID}{page Général}


\part{Référence}
\chapter{Tous les réglages}\label{chap:all-settings}
Ce chapitre est produit à partir de l'application elle-même : il ne peut donc pas être en retard sur elle. Vous y
trouvez chaque réglage dans l'ordre de l'écran des réglages, avec ses choix possibles, sa valeur par défaut et sa
clé interne. Cette clé vous servira surtout si vous nous signalez un problème : elle désigne le réglage sans
ambiguïté.
\capturedouble{prefs-en-top}{prefs-en-3}{Réglages, début}{suite}
% generated by tools/settings.py, do not edit
\section{Réglages}
\subsection*{Autres}
\begin{description}
\item[Durée par défaut]  \emph{Choix :} 30 secondes, 1 minute, 2 minutes, 3 minutes, 5 minutes, 10 minutes, 15 minutes, 30 minutes. \emph{Par défaut :} 3 minutes. {\footnotesize\texttt{zxtune.core.plugins.default\_duration}}
\item[Silence detection]  \emph{Choix :} Aucun, 3 secondes, 4 secondes, 5 secondes, 10 secondes, 15 secondes. \emph{Par défaut :} Aucun. {\footnotesize\texttt{zxtune.sound.silencelimit}}
\item[Loops limit]  \emph{Choix :} None, 1 loop, 2 loops, 3 loops, 4 loops, 5 loops. \emph{Par défaut :} None. {\footnotesize\texttt{zxtune.sound.looplimit}}
\item[Fade-in]  \emph{Choix :} Aucun, 1 seconde, 2 secondes, 3 secondes, 5 secondes, 10 secondes. \emph{Par défaut :} Aucun. {\footnotesize\texttt{zxtune.sound.fadein}}
\item[Fade-out]  \emph{Choix :} Aucun, 1 seconde, 2 secondes, 3 secondes, 5 secondes, 10 secondes. \emph{Par défaut :} Aucun. {\footnotesize\texttt{zxtune.sound.fadeout}}
\item[Gain] (plage : 50\% à 200\%) \emph{Par défaut :} 100\%. {\footnotesize\texttt{zxtune.sound.gain}}
\end{description}
\subsection*{AY–3–8910/YM2149F}
\begin{description}
\item[Fréquence d'horloge]  \emph{Choix :} ZX Spectrum (1,7734 MHz), Pentagon (1,75 MHz), Taganrog (3,5 MHz), Atari ST (2 MHz), Amstrad CPC (1 MHz). \emph{Par défaut :} Pentagon (1,75 MHz). {\footnotesize\texttt{zxtune.core.aym.clockrate}}
\item[Type]  \emph{Choix :} AY–3–8910, YM2149F. \emph{Par défaut :} AY–3–8910. {\footnotesize\texttt{zxtune.core.aym.type}}
\item[Interpolation]  \emph{Choix :} Aucune, Performance, Qualité. \emph{Par défaut :} Qualité. {\footnotesize\texttt{zxtune.core.aym.interpolation}}
\item[Disposition des canaux]  \emph{Choix :} ABC, ACB, BAC, BCA, CAB, CBA, Mono. \emph{Par défaut :} ABC. {\footnotesize\texttt{zxtune.core.aym.layout}}
\item[Canaux à rapport cyclique]  \emph{Choix :} Aucun, A, B, C, A, B, A, C, B, C, Tous. \emph{Par défaut :} Aucun. {\footnotesize\texttt{zxtune.core.aym.duty\_cycle\_mask}}
\item[Rapport cyclique] Largeur d'impulsion du son pour les canaux choisis, 50 = onde carrée (plage : 1\% à 99\%) \emph{Par défaut :} 50\%. {\footnotesize\texttt{zxtune.core.aym.duty\_cycle}}
\item[3-channel mixer]  {\footnotesize\texttt{zxtune.sound.mixer.3}}
\end{description}
\subsection*{Processeur Z80 (fichiers AY avec code du lecteur)}
\begin{description}
\item[Fréquence d'horloge]  \emph{Choix :} ZX Spectrum 48K, Pentagon (3,5 MHz), ZX Spectrum 128K (3,5469 MHz), Amstrad CPC (4 MHz), Turbo (7 MHz). \emph{Par défaut :} ZX Spectrum 48K, Pentagon (3,5 MHz). {\footnotesize\texttt{zxtune.core.z80.clockrate}}
\item[Durée d'interruption] Durée d'activité de la demande d'interruption à chaque image (plage : 1 T à 100 T) \emph{Par défaut :} 24 T. {\footnotesize\texttt{zxtune.core.z80.int\_ticks}}
\end{description}
\subsection*{YM2203 (TurboFM)}
\begin{description}
\item[Fréquence d'horloge]  \emph{Choix :} 1,75 MHz, TurboFM (3,5 MHz), NTSC (3,58 MHz), 4 MHz. \emph{Par défaut :} TurboFM (3,5 MHz). {\footnotesize\texttt{zxtune.core.fm.clockrate}}
\end{description}
\subsection*{YM2612/YM3438 (OPN2)}
\begin{description}
\item[Cœur d'émulation]  \emph{Choix :} Nuked (précis, CPU élevé), ymfm (MAME), GPGX (rapide). \emph{Par défaut :} Nuked (précis, CPU élevé). {\footnotesize\texttt{zxtune.core.fm.opn2.core}}
\item[Plateforme]  \emph{Choix :} Mega Drive modèle 1 (VA0-VA6, YM2612), Mega Drive modèle 1 VA7 (YM3438), Mega Drive modèle 2 (YM3438), Genesis 3 (YM3438), YM3438 sans filtre, YMF276 sans filtre, Personnalisé. \emph{Par défaut :} Mega Drive modèle 1 (VA0-VA6, YM2612). {\footnotesize\texttt{zxtune.core.fm.opn2.preset}}
\item[Puce]  \emph{Choix :} YM2612, YM3438, YMF276. \emph{Par défaut :} YM2612. {\footnotesize\texttt{zxtune.core.fm.opn2.chip}}
\item[Effet d'escalier du DAC] Distorsion de croisement du YM2612 discret, audible sur les sons faibles (cœurs Nuked et ymfm) \emph{Par défaut :} activé. {\footnotesize\texttt{zxtune.core.fm.opn2.ladder}}
\item[Filtre passe-bas] Coupure du filtre de la carte, 0 le désactive (plage : 0 Hz à 24000 Hz) \emph{Par défaut :} 3390 Hz. {\footnotesize\texttt{zxtune.core.fm.opn2.lowpass}}
\item[Pente du passe-bas]  \emph{Choix :} 6 dB/oct, 12 dB/oct. \emph{Par défaut :} 6 dB/oct. {\footnotesize\texttt{zxtune.core.fm.opn2.lowpass\_order}}
\item[Blocage du continu] Coupure des condensateurs de sortie, 0 garde la composante continue de la puce (plage : 0 Hz à 200 Hz) \emph{Par défaut :} 5 Hz. {\footnotesize\texttt{zxtune.core.fm.opn2.highpass}}
\item[Largeur stéréo] (plage : 0\% à 100\%) \emph{Par défaut :} 100\%. {\footnotesize\texttt{zxtune.core.fm.opn2.stereo}}
\item[Réponse de sortie]  {\footnotesize\texttt{ui.fm.opn2.response}}
\end{description}
\subsection*{YM2151 (OPM)}
\begin{description}
\item[Cœur d'émulation]  \emph{Choix :} Nuked (précis, CPU très élevé), ymfm (MAME), MAME (ancien, rapide). \emph{Par défaut :} ymfm (MAME). {\footnotesize\texttt{zxtune.core.fm.opm.core}}
\item[Plateforme]  \emph{Choix :} YM3012 (DAC stéréo : X68000, arcade), YM3014 (DAC mono), Personnalisé. \emph{Par défaut :} YM3012 (DAC stéréo : X68000, arcade). {\footnotesize\texttt{zxtune.core.fm.opm.preset}}
\item[Filtre passe-bas] Coupure du filtre de la carte, 0 le désactive (plage : 0 Hz à 24000 Hz) \emph{Par défaut :} 0 Hz. {\footnotesize\texttt{zxtune.core.fm.opm.lowpass}}
\item[Pente du passe-bas]  \emph{Choix :} 6 dB/oct, 12 dB/oct. \emph{Par défaut :} 12 dB/oct. {\footnotesize\texttt{zxtune.core.fm.opm.lowpass\_order}}
\item[Blocage du continu] Coupure des condensateurs de sortie, 0 garde la composante continue de la puce (plage : 0 Hz à 200 Hz) \emph{Par défaut :} 0 Hz. {\footnotesize\texttt{zxtune.core.fm.opm.highpass}}
\item[Largeur stéréo] (plage : 0\% à 100\%) \emph{Par défaut :} 100\%. {\footnotesize\texttt{zxtune.core.fm.opm.stereo}}
\item[Réponse de sortie]  {\footnotesize\texttt{ui.fm.opm.response}}
\end{description}
\subsection*{YM2203/YM2608/YM2610 (OPN)}
\begin{description}
\item[Cœur d'émulation]  \emph{Choix :} ymfm (MAME), MAME (ancien). \emph{Par défaut :} ymfm (MAME). {\footnotesize\texttt{zxtune.core.fm.opn.core}}
\item[Plateforme]  \emph{Choix :} PC-88/PC-98, Personnalisé. \emph{Par défaut :} PC-88/PC-98. {\footnotesize\texttt{zxtune.core.fm.opn.preset}}
\item[Filtre passe-bas] Coupure du filtre de la carte, 0 le désactive (plage : 0 Hz à 24000 Hz) \emph{Par défaut :} 0 Hz. {\footnotesize\texttt{zxtune.core.fm.opn.lowpass}}
\item[Pente du passe-bas]  \emph{Choix :} 6 dB/oct, 12 dB/oct. \emph{Par défaut :} 12 dB/oct. {\footnotesize\texttt{zxtune.core.fm.opn.lowpass\_order}}
\item[Blocage du continu] Coupure des condensateurs de sortie, 0 garde la composante continue de la puce (plage : 0 Hz à 200 Hz) \emph{Par défaut :} 0 Hz. {\footnotesize\texttt{zxtune.core.fm.opn.highpass}}
\item[Largeur stéréo] (plage : 0\% à 100\%) \emph{Par défaut :} 100\%. {\footnotesize\texttt{zxtune.core.fm.opn.stereo}}
\item[Réponse de sortie]  {\footnotesize\texttt{ui.fm.opn.response}}
\end{description}
\subsection*{YM3812/YMF262 (OPL)}
\begin{description}
\item[Cœur d'émulation]  \emph{Choix :} Nuked (précis), AdLibEmu (rapide), MAME (rapide). \emph{Par défaut :} Nuked (précis). {\footnotesize\texttt{zxtune.core.fm.opl.core}}
\item[Plateforme]  \emph{Choix :} AdLib (mono), Sound Blaster 1.x/2.0, Sound Blaster Pro 1/2, Sound Blaster 16, Personnalisé. \emph{Par défaut :} Sound Blaster 16. {\footnotesize\texttt{zxtune.core.fm.opl.preset}}
\item[Filtre passe-bas] Coupure du filtre de la carte, 0 le désactive (plage : 0 Hz à 24000 Hz) \emph{Par défaut :} 0 Hz. {\footnotesize\texttt{zxtune.core.fm.opl.lowpass}}
\item[Pente du passe-bas]  \emph{Choix :} 6 dB/oct, 12 dB/oct. \emph{Par défaut :} 6 dB/oct. {\footnotesize\texttt{zxtune.core.fm.opl.lowpass\_order}}
\item[Blocage du continu] Coupure des condensateurs de sortie, 0 garde la composante continue de la puce (plage : 0 Hz à 200 Hz) \emph{Par défaut :} 0 Hz. {\footnotesize\texttt{zxtune.core.fm.opl.highpass}}
\item[Largeur stéréo] (plage : 0\% à 100\%) \emph{Par défaut :} 100\%. {\footnotesize\texttt{zxtune.core.fm.opl.stereo}}
\item[Réponse de sortie]  {\footnotesize\texttt{ui.fm.opl.response}}
\end{description}
\subsection*{YM2413 (OPLL)}
\begin{description}
\item[Cœur d'émulation]  \emph{Choix :} Nuked (précis), EMU2413 (rapide), MAME (rapide). \emph{Par défaut :} Nuked (précis). {\footnotesize\texttt{zxtune.core.fm.opll.core}}
\item[Filtre passe-bas] Coupure du filtre de la carte, 0 le désactive (plage : 0 Hz à 24000 Hz) \emph{Par défaut :} 0 Hz. {\footnotesize\texttt{zxtune.core.fm.opll.lowpass}}
\item[Pente du passe-bas]  \emph{Choix :} 6 dB/oct, 12 dB/oct. \emph{Par défaut :} 12 dB/oct. {\footnotesize\texttt{zxtune.core.fm.opll.lowpass\_order}}
\item[Blocage du continu] Coupure des condensateurs de sortie, 0 garde la composante continue de la puce (plage : 0 Hz à 200 Hz) \emph{Par défaut :} 0 Hz. {\footnotesize\texttt{zxtune.core.fm.opll.highpass}}
\item[Réponse de sortie]  {\footnotesize\texttt{ui.fm.opll.response}}
\item[Interpolation]  \emph{Par défaut :} désactivé. {\footnotesize\texttt{zxtune.core.dac.interpolation}}
\item[Fréquence d'horloge]  \emph{Choix :} SAM Coupé (8 MHz), Creative Music System (7,16 MHz). \emph{Par défaut :} SAM Coupé (8 MHz). {\footnotesize\texttt{zxtune.core.saa.clockrate}}
\item[Interpolation]  \emph{Choix :} Aucune, Performance, Qualité. \emph{Par défaut :} Performance. {\footnotesize\texttt{zxtune.core.saa.interpolation}}
\end{description}
\subsection*{Écran principal}
\begin{description}
\item[Afficher les informations du morceau] Titre, puces, format du flux, émulation des puces, pilote d'origine et moteur au-dessus des commandes. Aussi basculé par le bouton à gauche des commandes \emph{Par défaut :} désactivé. {\footnotesize\texttt{ui.brief.expanded}}
\item[Afficher l'auteur] Auteur affiché sous le titre du morceau \emph{Par défaut :} désactivé. {\footnotesize\texttt{ui.brief.show\_author}}
\end{description}
\subsection*{Oscilloscope}
\begin{description}
\item[Oscilloscope] Affichage, déclenchement et performances {\footnotesize\texttt{ui.oscilloscope}}
\end{description}
\subsection*{Roland MT-32 (MIDI)}
\begin{description}
\item[Dossier des ROM] Choisir le dossier des dumps de ROM de contrôle et PCM du MT-32 ou CM-32L (entiers ou moitiés de puces, détectés par leur contenu) {\footnotesize\texttt{zxtune.core.mt32.roms\_path}}
\item[Modèle préféré]  \emph{Choix :} Automatique (MT-32 le plus récent d'abord), MT-32, CM-32L / LAPC-I. \emph{Par défaut :} Automatique (MT-32 le plus récent d'abord). {\footnotesize\texttt{zxtune.core.mt32.model}}
\item[Voix de l'oscilloscope]  \emph{Choix :} 9 parties (8 mélodiques + rythme), 32 partiels (synthèse LA). \emph{Par défaut :} 9 parties (8 mélodiques + rythme). {\footnotesize\texttt{zxtune.core.mt32.voices}}
\end{description}
\subsection*{SNES (SPC, SNSF)}
\begin{description}
\item[Emulation engine] ares: cycle accurate, plays stereo as the console. snes\_spc: lighter, without phase inverted surround. Applied from the next track \emph{Choix :} snes\_spc, ares. \emph{Par défaut :} ares. {\footnotesize\texttt{zxtune.core.spc.engine}}
\end{description}
\subsection*{MOS6581/SID}
\begin{description}
\item[Moteur d'émulation]  \emph{Choix :} reSID, reSIDfp. \emph{Par défaut :} reSIDfp. {\footnotesize\texttt{zxtune.core.sid.engine}}
\item[Modèle de puce par défaut] Utilisé si le morceau ne le précise pas \emph{Choix :} MOS6581, MOS8580. \emph{Par défaut :} MOS6581. {\footnotesize\texttt{zxtune.core.sid.model}}
\item[Forcer le modèle de puce] Ignorer la valeur indiquée dans le morceau \emph{Par défaut :} désactivé. {\footnotesize\texttt{zxtune.core.sid.model\_force}}
\item[Modèle de C64 par défaut] \%s (utilisé si le morceau ne le précise pas) \emph{Choix :} PAL, NTSC, NTSC (old), Drean (PAL-N), PAL-M. \emph{Par défaut :} PAL. {\footnotesize\texttt{zxtune.core.sid.clock}}
\item[Forcer le modèle de C64] Ignorer la valeur indiquée dans le morceau \emph{Par défaut :} désactivé. {\footnotesize\texttt{zxtune.core.sid.clock\_force}}
\item[Interpolation]  \emph{Choix :} Aucune, Performance, Qualité. \emph{Par défaut :} Aucune. {\footnotesize\texttt{zxtune.core.sid.interpolation}}
\item[zxtune.core.sid.filter]  \emph{Par défaut :} désactivé. {\footnotesize\texttt{zxtune.core.sid.filter}}
\item[Préréglage du filtre MOS6581]  \emph{Choix :} Custom, JSIDPlay2 reSID: Lightest, JSIDPlay2 reSID: Lighter, JSIDPlay2 reSID: Light, JSIDPlay2 reSID: Average, JSIDPlay2 reSID: Dark, JSIDPlay2 reSID: Darker, JSIDPlay2 reSID: Darkest, VICE default (reSID bias +500 mV), ReSID 6581 (resid-fp fit), Alankila 6581R4AR 3789, Alankila 6581R3 3984 \#1, Alankila 6581R3 3984 \#2, Lord Nightmare 6581R3 4285, Lord Nightmare 6581R3 4485, Lord Nightmare 6581R4 1986S, ZrX 6581R3 0384, ZrX 6581R3 1984, ZrX 6581R3 3684, ZrX 6581R3 3985, ZrX 6581R4AR 2286, Trurl 6581R3 0784, Trurl 6581R3 0486S, Trurl 6581R3 3384, Trurl 6581R3 4885, Trurl 6581R4AR 3789, Trurl 6581R4AR 4486, Nata 6581R3 2083, Grue 6581R4AR 3488, KruLLo, Enigma 6581R3 4885, Enigma 6581R3 1585. \emph{Par défaut :} Custom. {\footnotesize\texttt{zxtune.core.sid.filter\_6581\_preset}}
\item[Préréglage du filtre MOS8580]  \emph{Choix :} Custom, JSIDPlay2 reSID: Light, JSIDPlay2 reSID: Average, JSIDPlay2 reSID: Dark, Trurl 8580R5 1489, Trurl 8580R5 3691. \emph{Par défaut :} Custom. {\footnotesize\texttt{zxtune.core.sid.filter\_8580\_preset}}
\item[Courbe de coupure du filtre]  {\footnotesize\texttt{ui.sid.filter\_curve}}
\item[Digi boost du MOS8580] MOS8580 seulement : rend audibles les échantillons joués par le registre de volume \emph{Par défaut :} désactivé. {\footnotesize\texttt{zxtune.core.sid.digiboost}}
\item[reSID : biais du filtre MOS6581] Moteur reSID, filtre MOS6581 personnalisé seulement : décale la correspondance de la fréquence de coupure (plage : -1000 mV à 1000 mV) \emph{Par défaut :} 0 mV. {\footnotesize\texttt{zxtune.core.sid.filter\_bias}}
\item[reSIDfp : courbe du filtre MOS6581] Moteur reSIDfp, filtre personnalisé seulement : 0 = clair, 100 = sombre \emph{Par défaut :} 50. {\footnotesize\texttt{zxtune.core.sid.filter\_6581\_curve}}
\item[reSIDfp : plage du filtre MOS6581] Moteur reSIDfp, filtre personnalisé seulement : 0 = sombre, 100 = clair \emph{Par défaut :} 50. {\footnotesize\texttt{zxtune.core.sid.filter\_6581\_range}}
\item[reSIDfp : courbe du filtre MOS8580] Moteur reSIDfp, filtre personnalisé seulement : 0 = clair, 100 = sombre \emph{Par défaut :} 50. {\footnotesize\texttt{zxtune.core.sid.filter\_8580\_curve}}
\item[reSIDfp : formes d'onde combinées]  \emph{Choix :} Moyennes, Faibles, Fortes. \emph{Par défaut :} Moyennes. {\footnotesize\texttt{zxtune.core.sid.combined\_waveforms}}
\item[Impulsion+dent de scie du 6581 (reSIDfp)]  \emph{Choix :} Normale (la plupart des 6581), Forte (puces à impulsion+dent de scie forte). \emph{Par défaut :} Normale (la plupart des 6581). {\footnotesize\texttt{zxtune.core.sid.strong\_pulse\_saw}}
\item[Sortie 2SID/3SID]  \emph{Choix :} Stéréo (puces réparties), Mono. \emph{Par défaut :} Stéréo (puces réparties). {\footnotesize\texttt{zxtune.core.sid.multi\_sid\_mono}}
\item[Voie des samples (\$D418)] Oscilloscope : une voie en plus avec les samples du volume principal (digis 4 bits), les trois voies montrent les oscillateurs seuls. Appliqué au morceau suivant \emph{Par défaut :} désactivé. {\footnotesize\texttt{zxtune.core.sid.digi\_lane}}
\end{description}
\subsection*{Android Auto}
\begin{description}
\item[Android Auto] Options de l'écran de la voiture {\footnotesize\texttt{ui.car}}
\end{description}
\subsection*{Sauvegarde de la playlist}
\begin{description}
\item[Sauvegarde de la playlist] Exporter ou importer la playlist, exporter un journal de diagnostic {\footnotesize\texttt{rescue}}
\end{description}
\section{Oscilloscope}
\subsection*{Affichage}
\begin{description}
\item[Fenêtre des formes d'onde] Durée affichée de la forme d'onde de chaque voix (plage : 5 ms à 100 ms) \emph{Par défaut :} 40 ms. {\footnotesize\texttt{ui.oscilloscope.window\_ms}}
\item[Amplification visuelle]  \emph{Choix :} Automatique (par voix), $\times$1, $\times$1,5, $\times$2, $\times$3, $\times$4, $\times$6, $\times$8. \emph{Par défaut :} Automatique (par voix). {\footnotesize\texttt{ui.oscilloscope.gain}}
\item[Périodes affichées au plus (zoom sur les notes aiguës)]  \emph{Choix :} Désactivé (toute la fenêtre), 4 périodes, 8 périodes, 16 périodes, 32 périodes. \emph{Par défaut :} 8 périodes. {\footnotesize\texttt{zxtune.sound.scope.max\_periods}}
\item[Fenêtre des jauges SID] Durée affichée des jauges des oscillateurs et du volume général (plage : 2 ms à 40 ms) \emph{Par défaut :} 10 ms. {\footnotesize\texttt{ui.oscilloscope.gauge\_window\_ms}}
\item[Hauteur de l'oscilloscope] Par rapport au 16:9 de la largeur, le tableau de bord en dessous prend le reste de l'écran (plage : 60\% à 160\%) \emph{Par défaut :} 100\%. {\footnotesize\texttt{ui.oscilloscope.height}}
\item[Hauteur de la réponse du filtre] Part de l'oscilloscope prise par la réponse du filtre SID au-dessus des formes d'onde des voix (plage : 20\% à 80\%) \emph{Par défaut :} 55\%. {\footnotesize\texttt{ui.oscilloscope.filter\_share}}
\item[Voies de batterie AdLib] OPL/OPL2 : voies des batteries du mode rythme (peu de musiques l'utilisent), sinon 9 voix en 3$\times$3. Appliqué au morceau suivant \emph{Par défaut :} désactivé. {\footnotesize\texttt{ui.oscilloscope.opl\_rhythm}}
\item[Formes d'onde sous la réponse du filtre] La vue du filtre SID montre la courbe au-dessus des formes d'onde des voix, sinon la courbe seule \emph{Par défaut :} activé. {\footnotesize\texttt{ui.oscilloscope.filter\_waves}}
\item[Afficher les statistiques techniques] Lignes de l'oscilloscope, de la sortie son, du rendu et du GPU sous le tableau de bord \emph{Par défaut :} activé. {\footnotesize\texttt{ui.oscilloscope.show\_stats}}
\item[Limite d'images par seconde]  \emph{Choix :} Fréquence de l'écran, 120 i/s, 90 i/s, 60 i/s, 30 i/s. \emph{Par défaut :} Fréquence de l'écran. {\footnotesize\texttt{ui.oscilloscope.max\_fps}}
\item[Moteur graphique]  \emph{Choix :} Automatique (Vulkan si disponible), Vulkan, OpenGL ES. \emph{Par défaut :} Automatique (Vulkan si disponible). {\footnotesize\texttt{ui.oscilloscope.backend}}
\item[Pilote Vulkan personnalisé] Choisir le paquet du pilote (zip avec meta.json ou bibliothèque .so, p. ex. Mesa Turnip pour les GPU Adreno) {\footnotesize\texttt{ui.oscilloscope.vulkan\_driver}}
\item[Utiliser le pilote Vulkan personnalisé] Le pilote du système est utilisé si le pilote personnalisé ne se charge pas \emph{Par défaut :} désactivé. {\footnotesize\texttt{ui.oscilloscope.vulkan\_driver.enabled}}
\end{description}
\subsection*{Déclenchement}
\begin{description}
\item[Préréglage]  \emph{Choix :} Par défaut (modèle corrscope, 40 ms), Notes évolutives (60 ms), Chiptune carré/impulsion, Basses, notes graves, Triangle NES, Économie d'énergie, Personnalisé. \emph{Par défaut :} Par défaut (modèle corrscope, 40 ms). {\footnotesize\texttt{ui.oscilloscope.trigger\_preset}}
\item[Largeur de déclenchement] Taille de l'historique, doit dépasser deux périodes des notes les plus graves (plage : 10 ms à 200 ms) \emph{Par défaut :} 40 ms. {\footnotesize\texttt{zxtune.sound.scope.trigger.width}}
\item[Force des fronts] Préférer les fronts montants raides : les nouvelles notes sont mieux saisies, mais les notes tenues sautent davantage (plage : 0\% à 300\%) \emph{Par défaut :} 100\%. {\footnotesize\texttt{zxtune.sound.scope.trigger.edge\_strength}}
\item[Largeur de pente] (plage : 5 \% de la période à 100 \% de la période) \emph{Par défaut :} 25 \% de la période. {\footnotesize\texttt{zxtune.sound.scope.trigger.slope\_width}}
\item[Force de l'historique] Préférer la ressemblance avec les images précédentes (plage : 0\% à 300\%) \emph{Par défaut :} 100\%. {\footnotesize\texttt{zxtune.sound.scope.trigger.buffer\_strength}}
\item[Réactivité de l'historique] (plage : 0\% à 100\%) \emph{Par défaut :} 50\%. {\footnotesize\texttt{zxtune.sound.scope.trigger.responsiveness}}
\item[Atténuation de l'historique] (plage : 10 \% de la période à 200 \% de la période) \emph{Par défaut :} 50 \% de la période. {\footnotesize\texttt{zxtune.sound.scope.trigger.buffer\_falloff}}
\item[Réinitialiser sous la ressemblance] Oublier l'historique aux nouvelles notes si l'onde ne lui ressemble pas (plage : 0\% à 100\%) \emph{Par défaut :} 30\%. {\footnotesize\texttt{zxtune.sound.scope.trigger.reset\_below}}
\item[Vitesse de suppression du continu] (plage : 0\% à 100\%) \emph{Par défaut :} 0\%. {\footnotesize\texttt{zxtune.sound.scope.trigger.mean\_responsiveness}}
\item[Amplification du signe] Passages par zéro raides pour les ondes sans fronts nets (p. ex. triangle) (plage : 0\% à 300\%) \emph{Par défaut :} 0\%. {\footnotesize\texttt{zxtune.sound.scope.trigger.sign\_strength}}
\item[Plage de recherche] (plage : 10 \% de la largeur à 100 \% de la largeur) \emph{Par défaut :} 50 \% de la largeur. {\footnotesize\texttt{zxtune.sound.scope.trigger.diameter}}
\item[Fronts descendants] P. ex. pour les trompettes échantillonnées \emph{Par défaut :} désactivé. {\footnotesize\texttt{zxtune.sound.scope.trigger.falling}}
\item[Suivi de hauteur] Remettre l'historique à l'échelle quand la hauteur change \emph{Par défaut :} activé. {\footnotesize\texttt{zxtune.sound.scope.trigger.pitch\_tracking}}
\item[Recalage au passage par zéro, échantillons] Glisser jusqu'au passage par zéro montant le plus proche, 0 pour désactiver \emph{Par défaut :} 0. {\footnotesize\texttt{zxtune.sound.scope.trigger.post\_radius}}
\item[Recherches complètes]  \emph{Choix :} À chaque image, 60 par seconde, 30 par seconde. \emph{Par défaut :} À chaque image. {\footnotesize\texttt{zxtune.sound.scope.trigger.search\_rate}}
\item[Sous-échantillonnage du déclenchement]  \emph{Choix :} Automatique, Aucun (pleine fréquence), 2, 4. \emph{Par défaut :} Automatique. {\footnotesize\texttt{zxtune.sound.scope.trigger.subsampling}}
\end{description}
\section{Mixage}
\subsection*{Channel A}
\begin{description}
\item[Left] (plage : 0\% à 100\%) \emph{Par défaut :} 100\%. {\footnotesize\texttt{zxtune.sound.mixer.3.0\_0}}
\item[Right] (plage : 0\% à 100\%) \emph{Par défaut :} 0\%. {\footnotesize\texttt{zxtune.sound.mixer.3.0\_1}}
\end{description}
\subsection*{Channel B}
\begin{description}
\item[Left] (plage : 0\% à 100\%) \emph{Par défaut :} 60\%. {\footnotesize\texttt{zxtune.sound.mixer.3.1\_0}}
\item[Right] (plage : 0\% à 100\%) \emph{Par défaut :} 60\%. {\footnotesize\texttt{zxtune.sound.mixer.3.1\_1}}
\end{description}
\subsection*{Channel C}
\begin{description}
\item[Left] (plage : 0\% à 100\%) \emph{Par défaut :} 0\%. {\footnotesize\texttt{zxtune.sound.mixer.3.2\_0}}
\item[Right] (plage : 0\% à 100\%) \emph{Par défaut :} 100\%. {\footnotesize\texttt{zxtune.sound.mixer.3.2\_1}}
\end{description}
\section{Android Auto}
\subsection*{Android Auto}
\begin{description}
\item[Afficher l'oscilloscope / le tableau de bord dans la voiture] Remplace la pochette sur l'écran de la voiture par une visualisation animée pendant la lecture. Lisez l'avertissement légal avant d'activer. \emph{Par défaut :} désactivé. {\footnotesize\texttt{ui.car.visualization}}
\item[Contenu affiché dans la voiture]  \emph{Choix :} Oscilloscope, Tableau de bord des puces, Les deux. \emph{Par défaut :} Oscilloscope. {\footnotesize\texttt{ui.car.visualization.mode}}
\item[Fréquence d'images dans la voiture]  \emph{Choix :} 1 img/s, 2 img/s, 5 img/s, 10 img/s, 15 img/s. \emph{Par défaut :} 5 img/s. {\footnotesize\texttt{ui.car.visualization.fps}}
\item[Avertissement légal] Règles sur les écrans visibles par le conducteur dans votre pays {\footnotesize\texttt{ui.car.legal}}
\end{description}


\chapter{Statistiques techniques}
Sous le tableau de bord, quelques lignes techniques décrivent ce que fait votre téléphone. Vous n'en aurez besoin
qu'en cas de souci, mais elles disent vite où est le problème.
\begin{description}
  \item[scope] le nombre de voies affichées, le coût du déclenchement et l'écart de synchronisation entre le son et
        l'image.
  \item[output] la fréquence de sortie et la taille du tampon audio. Un tampon plus grand protège des coupures au
        prix d'un léger retard de l'image sur le son.
  \item[render] la vitesse de l'émulation par rapport au temps réel. C'est la ligne la plus utile : à 300~\%, votre
        téléphone calcule trois secondes de musique par seconde ; sous 100~\%, il ne suit plus et le son coupe.
  \item[vulkan / gles] le moteur de dessin, les images par seconde et le temps de dessin de chacune.
  \item[vk] le processeur graphique et la version de son pilote.
\end{description}

\chapter{Dépannage}
\begin{description}
  \item[Craquements] baisser la fréquence d'images de l'oscilloscope, choisir un cœur d'émulation plus léger, ou
        augmenter le tampon audio.
  \item[Fichier non reconnu] vérifier qu'il figure dans la table des formats ; certains formats demandent des
        fichiers annexes (bibliothèques \_lib des rips xSF) dans le même dossier.
  \item[Catalogue indisponible] le service en ligne peut être hors ligne ; réessayer plus tard.
\end{description}
\arediger{Problèmes rencontrés par les utilisateurs et leurs solutions.}

\chapter{Limites connues}
Nous préférons vous dire ce qui ne marche pas encore plutôt que de vous laisser le découvrir.
\begin{itemize}
  \item Android Auto affiche au plus 32 voies : au-delà, les dernières ne sont pas montrées.
  \item Les fichiers Saturn SSF et PlayStation PSF n'ont pas encore de vue détaillée par voix ; leur oscilloscope
        montre la sortie générale.
  \item La vue PC Engine ne sert que pour les enregistrements VGM ; les fichiers HES utilisent la vue générique.
  \item Les réglages marqués « applied to the next track » ne changent pas le morceau en cours.
  \item L'interface est en anglais uniquement. Les traductions existantes sont conservées dans les sources mais ne
        sont plus livrées, faute de pouvoir les tenir à jour.
  \item Le MT-32 a besoin des images mémoire du module d'origine, que nous ne fournissons pas.
  \item Nous ne développons pour l'instant que la version Android ; la version de bureau suit avec du retard.
\end{itemize}


\appendix
\chapter{Glossaire}
\begin{description}
  \item[ADSR] attaque, déclin, maintien, relâchement : les quatre phases de l'enveloppe d'un son. L'attaque est la
        montée au début de la note, le déclin la descente vers le niveau de maintien, le relâchement l'extinction
        après la fin de la note.
  \item[Algorithme] sur une puce FM, la façon dont les opérateurs d'une voix sont branchés entre eux.
  \item[Arpège] accord joué en faisant défiler très vite ses notes sur une seule voix, pour donner l'illusion de
        plusieurs voix.
  \item[Cent] centième de demi-ton.
  \item[Coupure] fréquence à partir de laquelle un filtre commence à atténuer le son.
  \item[DDA] mode du HuC6280 où le processeur écrit directement les échantillons dans la voix.
  \item[Déclenchement] calage de la fenêtre de l'oscilloscope sur la période du son, pour que la forme reste
        immobile.
  \item[Driver] programme qui joue la musique sur la machine d'origine en écrivant dans les registres de la puce.
  \item[DSP] processeur de traitement du signal, utilisé pour l'écho et la réverbération.
  \item[Enveloppe] évolution du volume (ou d'un autre paramètre) d'une note dans le temps.
  \item[FIR] filtre à réponse impulsionnelle finie ; sur la Super Nintendo, il colore les répétitions de l'écho.
  \item[FM] synthèse par modulation de fréquence : un oscillateur fait varier la fréquence d'un autre, ce qui crée
        des harmoniques.
  \item[Modulation en anneau] multiplication de deux signaux, qui donne des sons métalliques ou de cloche.
  \item[Opérateur] oscillateur sinusoïdal d'une voix FM, avec sa propre enveloppe.
  \item[PAL-N] norme vidéo argentine : horloge proche du NTSC, 50 images par seconde.
  \item[Porteuse] opérateur FM que l'on entend ; les modulateurs ne s'entendent qu'à travers elle.
  \item[Rapport cyclique] proportion du temps où une onde carrée est en position haute ; il change le timbre.
  \item[Registre] case mémoire de la puce dans laquelle le driver écrit un réglage.
  \item[Résonance] accentuation des fréquences autour de la coupure d'un filtre.
  \item[Rip] fichier extrait d'un jeu, qui contient le code et les données de sa musique.
  \item[Synchronisation] remise à zéro d'un oscillateur à chaque cycle d'un autre, qui donne des sons durs et
        riches.
  \item[Table d'ondes] forme d'onde mémorisée, rejouée en boucle.
  \item[Tracker] logiciel de composition par motifs, et format des musiques qu'il produit.
  \item[Voie] ligne de l'oscilloscope qui montre une voix.
\end{description}

\chapter{Historique des versions}
Cette annexe résume les grandes étapes du développement de \soft, des plus récentes aux plus anciennes. Les
numéros sont ceux des versions de développement remises pour essai.

\section*{Version 40}
\begin{itemize}
  \item Interface entièrement en anglais.
  \item SID : la colonne Mod revient à la seule voix source (« Sync V1 », « Sync+Ring V1 »).
\end{itemize}
\section*{Version 39}
\begin{itemize}
  \item Vue PC Engine : forme de l'onde décrite en mots (sinus, triangle, scie montante ou descendante, carré,
        impulsion), numéro de table et amplitude.
\end{itemize}
\section*{Version 38}
\begin{itemize}
  \item Vue PC Engine : barre de niveau et barres de panoramique.
  \item SID : la voie des samples se place sous les voix de sa puce ; les trois voix ne contiennent plus les
        échantillons joués par le volume.
  \item AdLib OPL/OPL2 : neuf voix en 3~$\times$~3, voies de batterie masquées par défaut.
\end{itemize}
\section*{Version 37}
\begin{itemize}
  \item PC Engine : sortie par voix et état du HuC6280 dans le cœur Ootake ; vue dédiée avec la table d'ondes de
        chaque voix et tous ses registres.
\end{itemize}
\section*{Version 36}
\begin{itemize}
  \item Plein écran portrait : encoche du haut et coins du bas évités, tableau de bord correctement placé.
  \item SID : la colonne Mod indique la voix source de la synchronisation et de l'anneau.
\end{itemize}
\section*{Version 35}
\begin{itemize}
  \item Réglages de hauteur de l'oscilloscope et de la réponse du filtre.
  \item SID : voie des samples (\$D418) en option.
  \item Panneau en direct : page SNES.
\end{itemize}
\section*{Versions 33 et 34}
\begin{itemize}
  \item Listes de lecture multiples, favoris, ajout depuis le lecteur, le navigateur et les listes.
  \item File d'attente « Jouer ensuite », import et export M3U, bouton $\equiv$+ dans le lecteur.
  \item Fichiers multi-morceaux lus directement depuis le navigateur.
  \item SID : sortie mono en option pour les morceaux 2SID et 3SID.
\end{itemize}
\section*{Versions 30 à 32}
\begin{itemize}
  \item Super Nintendo : émulation ares (au cycle près) par défaut, snes\_spc en option, pour les SPC et les SNSF.
\end{itemize}
\section*{Versions antérieures}
\begin{itemize}
  \item Vue de la réponse du filtre SID dans l'oscilloscope, sans texte de valeurs.
  \item DSP d'effets du SCSP activé : réverbération des bornes Sega et de la Saturn.
  \item Valeurs du tableau de bord en temps réel, image par image.
  \item Informations techniques du morceau (codec, fréquence, résolution, débit), bas d'écran compact.
  \item Cœurs FM au choix par famille, préréglages de plateforme et étage de sortie analogique.
  \item Vue des puces PCM YMW258 et SCSP.
  \item WAV étendu (PCM et virgule flottante).
  \item Identité LHUNA : barre de navigation, mini-lecteur, réglages en cartes, À propos.
  \item Oscilloscope Vulkan et OpenGL ES à traits fins, fenêtre adaptée à la période de chaque voix.
  \item Inspecteur FM (OPN, OPM, OPL) et vue S-DSP.
\end{itemize}

\chapter{Crédits et licences des composants}\label{chap:credits}
\soft{} repose sur le travail de nombreux projets libres. La table ci-dessous est générée à partir des fichiers de
licence présents dans les sources ; le texte complet de chaque licence accompagne le code source et figure dans
l'écran « À propos » de l'application.
Base de l'application : \texttt{zxtune}, sous LGPL~v3.
% generated by tools/credits.py, do not edit
\begin{longtable}{@{}l l@{}}
\toprule \textbf{Composant (dossier)} & \textbf{Licence} \\ \midrule \endhead
\texttt{FLAC} & BSD (Xiph) \\
\texttt{adrenotools} & BSD 2-clause \\
\texttt{asap} & GPL v2 \\
\texttt{atrac9} & MIT \\
\texttt{curl} & curl (MIT-like) \\
\texttt{ffmpeg} & LGPL v2.1 \\
\texttt{gme/gme} & LGPL v2.1 \\
\texttt{he} & GPL v3 \\
\texttt{ht} & GPL v3 \\
\texttt{ht/Core} & GPL v3 \\
\texttt{hvl} & BSD 3-clause \\
\texttt{lame} & LGPL v2 \\
\texttt{lhasa} & ISC \\
\texttt{mgba} & MPL 2.0 \\
\texttt{minimp3} & CC0 \\
\texttt{mpg123} & LGPL v2.1 \\
\texttt{ogg} & BSD \\
\texttt{openmpt} & BSD \\
\texttt{opus} & BSD \\
\texttt{sidplayfp} & GPL v2 \\
\texttt{snesspc} & LGPL v2.1 \\
\texttt{snesspc/ares} & ISC \\
\texttt{sseqplayer} & WTFPL \\
\texttt{unrar} & UnRAR (freeware, source available) \\
\texttt{v2m} & Artistic 2.0 \\
\texttt{vgmstream} & ISC-like (see COPYING) \\
\texttt{vio2sf/desmume} & GPL v2 \\
\texttt{vorbis} & BSD \\
\texttt{xmp} & MIT \\
\texttt{z80ex} & GPL v2 \\
\bottomrule \end{longtable}

La bibliothèque libvgm (dossier \texttt{vgm}) et les cœurs qu'elle contient (Nuked, MAME, Ootake, EMU2413\dots)
n'ont pas de fichier de licence unique dans les sources : la licence de chaque cœur doit être vérifiée
individuellement avant publication de l'APK sous GPL~v3.

\capturedouble{about-credits-en}{about-licenses}{Crédits dans l'application}{licences}
\section*{Principaux auteurs}
\begin{description}
  \item[ZXTune] Vitamin, base de l'application.
  \item[libsidplayfp] Simon White, Antti Lankila, Leandro Nini et contributeurs ; reSID par Dag Lem.
  \item[ares] Near et l'équipe ares.
  \item[snes\_spc, Game Music Emu] Shay Green (blargg) et contributeurs.
  \item[libvgm] Valley Bell ; cœurs Nuked par Alexey Khokholov, ymfm par Aaron Giles, Ootake par Kitao Nakamura.
  \item[vgmstream] Adam Gashlin, bnnm et contributeurs.
  \item[libopenmpt] l'équipe OpenMPT ; \textbf{libxmp} Claudio Matsuoka et contributeurs.
  \item[Munt] l'équipe Munt.
  \item[ASAP] Piotr Fusik et contributeurs.
  \item[Highly Experimental / Theoretical] Neill Corlett (cœurs PlayStation et Saturn).
\end{description}
Les sites de ces projets sont listés en bibliographie.

\chapter{Bibliographie}

\section*{Projets logiciels}
\begin{itemize}
  \item ZXTune, par Vitamin : \url{https://github.com/vitamin-caig/zxtune}
  \item libsidplayfp, reSID et reSIDfp : \url{https://github.com/libsidplayfp/libsidplayfp}
  \item ares : \url{https://ares-emu.net}
  \item libvgm, par Valley Bell : \url{https://github.com/ValleyBell/libvgm}
  \item Cœurs Nuked, par nukeykt : \url{https://github.com/nukeykt}
  \item ymfm, par Aaron Giles : \url{https://github.com/aaronsgiles/ymfm}
  \item vgmstream : \url{https://github.com/vgmstream/vgmstream}
  \item libopenmpt : \url{https://lib.openmpt.org}
  \item Munt (émulateur MT-32) : \url{https://github.com/munt/munt}
  \item ASAP (POKEY) : \url{https://asap.sourceforge.net}
\end{itemize}

\section*{Collections}
\begin{itemize}
  \item High Voltage SID Collection : \url{https://www.hvsc.c64.org}
  \item The Mod Archive : \url{https://modarchive.org}
  \item VGMRips : \url{https://vgmrips.net}
  \item ZXArt : \url{https://zxart.ee}
\end{itemize}

\section*{Documentation des puces}
\begin{itemize}
  \item Fiches techniques des fabricants (General Instrument AY-3-8910, MOS 6581, Yamaha YM2612, YM3812, YMF262,
        YM2151, YMF292) et documentation des projets d'émulation ci-dessus.
  \item Documentation communautaire du matériel : NESdev Wiki (NES), fullsnes (Super Nintendo), Pan Docs
        (Game Boy), documentation de la scène PC Engine.
\end{itemize}

\section*{Références des réglages}
\begin{itemize}
  \item JSIDPlay2 (préréglages de filtre SID) : \url{https://sourceforge.net/projects/jsidplay2/}
  \item VICE (biais du filtre reSID) : \url{https://vice-emu.sourceforge.io}
  \item sidplayfp, page de manuel \texttt{sidplayfp.ini} : \url{https://www.mankier.com/5/sidplayfp.ini}
  \item DOSBox Staging, émulation des filtres de sortie : \url{https://www.dosbox-staging.org/0.83/manual/sound/output-filters/}
  \item BlastEm : \url{https://www.retrodev.com/blastem/}
  \item Genesis Plus GX, discussion du filtre passe-bas : \url{https://forums.libretro.com/t/genesis-plus-gx-where-is-the-audio-lowpass-filter/12169}
  \item Console Mods, circuit audio de la Genesis : \url{https://consolemods.org/wiki/Genesis:Audio_Circuit_Mod}
  \item Mesen, options audio SNES : \url{https://www.mesen.ca/snes/docs/configuration/audio.html}
  \item Snes9x (cœur libretro), interpolation : \url{https://docs.libretro.com/library/snes9x/}
  \item corrscope (déclencheur par corrélation) : \url{https://github.com/corrscope/corrscope}
\end{itemize}


\end{document}
